\documentclass{article}
\usepackage{packages/notespkg}
\addbibresource{refs.bib}




\title{Coalgebra, Categories, Automota, and MDPs}
\author{lorisj }
\date{May 2025}

\begin{document}

\maketitle

\newpage
\tableofcontents

\newpage
\section{Introduction}
The main goal of these notes is to explore new ideas in the intersection of:
\begin{itemize}
    \item Category theory/coalgebra
    \item MDPs (Markov Decision Processes)
    \item LDBA (Limit Deterministic Buchi Automota),
    \item Bisimulation and bisimulation metrics
\end{itemize}

and also serve as a reference for some results/definitions in these areas.

\newpage
\section{Notation}
\begin{defB}{Notation Used}
    These are default notations that are used unless otherwise noted:
    \begin{itemize}
        \item TFDC $\eqdef$ The following diagram(s) commute(s).
        \item s.t. $\eqdef$ such that
        \item $A^B \eqdef \{f : B \to A\}$ for sets $A,B$.
        \item $2^A = \mathcal P[A]$ (powerset of $A$) for set $A$, note that this means the powerset can be interpreted as the set of functions from $A \to 2$.
        \item $\omega\eqdef \{\text{natural numbers including 0} \} =\{0,1,2,\ldots \}$.
        \item  $\forall_{v}^V, \exists_v^V \eqdef \forall_{v \in V}, \exists_{v \in V}$, where the superscript refers to the set being quantified over (domain) and the subscript denotes the variable bound by the quantifier.
        \item $f[X], f^{-1}[X]\eqdef$ image, preimage for a set $X$ under function $f$ (in contrast $f(x),f^{-1}(x)$ denote element application)
        \item $(\cdot)$ denotes a unique free variable (i.e. $(\cdot) +2$ is the function that takes in a variable and adds 2)
        \item $\Delta X$ \refdef{distribution-functor}
    \end{itemize}
\end{defB}


\newpage
\section{Category Theory}

\subsection{Categories}
Intuitively, a category can be thought of as a more structured graph, where you can compose edges. Specifically, a category is a directed multigraph, and also there is a composition operation, as well as an identity. The nodes are called Objects, and the edges are called arrows (or morphisms). Essentially a category will look something like this:
\begin{center}
    % https://tikzcd.yichuanshen.de/#N4Igdg9gJgpgziAXAbVABwnAlgFyxMJZABgBpiBdUkANwEMAbAVxiRAEEQBfU9TXfIRQBGclVqMWbAELdeIDNjwEiogEzj6zVohAAROXyWCiAZjHUtU3QGFu4mFADm8IqABmAJwgBbJGuocCCRhHg9vP0RRECCkYjCQL19-QODEMhAGLDAdEDgILKhDRIi41LLMiAg0IjUAdjJ3RjgYcQY6ACMYBgAFfmUhEE8sJwALHGKkyOjYqOoGKpqUAE5G5tb5zu6+4xVdBhh3CYSplJi0gMrqojWGFrat3v6TXWGx4-lTxHNzpB+F64oW73TZdJ67QZvcaTUrfcpwq5LZCrUhNO4bTKPHYCPZDEbQk6wy6zDJdMBFRAANniFC4QA
    \begin{tikzcd}
        A \arrow[loop, distance=2em, in=215, out=145] & B \arrow[l] \arrow[loop, distance=2em, in=35, out=325]                                         &  & C \arrow[loop, distance=2em, in=305, out=235] \arrow[loop, distance=2em, in=35, out=325] \\
        &                                                                                                &  &                                                                                          \\
        & D \arrow[uu] \arrow[luu] \arrow[loop, distance=2em, in=305, out=235] \arrow[luu, bend left=60] &  &
    \end{tikzcd}
\end{center}

The important notes are that:
\begin{enumerate}
    \item You can have multiple arrows between two nodes that are not equal
    \item If I have two arrows: $f: A \to B, g:B \to C$, then there must exist some arrow $g \circ f : A \to C$.
    \item There exists at least one self loop for each object, called the identity, $\text{id}_c : C \to C$ such that its composition with any other object does not change it.
\end{enumerate}




More formally:
\begin{defB}[category]{Category}
    A category $C$ consists of:
    \begin{itemize}
        \item $\textbf{Obj}(C)$: A class of objects
        \item $\Hom(C)$: A class of morphisms (sometimes called arrows)
        \item $ \circ$: A class function $: \Hom(C) \times \Hom(C) \to \Hom(C)$
    \end{itemize}
    Such that:

    \begin{enumerate}
        \item Each morphism has a \emph{domain} and \emph{codomain}:

              For every $f \in \Hom(C)$ there must be two objects, a \emph{domain} and \emph{codomain}, both in $\mathbf{Obj}(C)$.


              As notation, $f:X\to Y$ denotes a morphism with domain $X$, codomain $Y$, and $\Hom (X,Y)$ for objects $X,Y$ denotes the class of all morphisms with domain $X$, codomain $Y$.

        \item Composition is \emph{type-respecting}:

              If $f: X \to Y$ and $g: Y \to Z$, then the composite $g \circ f$ must exist in $\Hom C$, and must be of type: $g \circ f : X \to Z$.

        \item Composition is \emph{associative}:

              For morphisms $f: W \to X$, $g: X \to Y$, and $h: Y \to Z$:
              \begin{align}
                  (h \circ g) \circ f \;=\; h \circ (g \circ f)
              \end{align}

        \item \emph{Identity morphisms} exist:


              Every $X \in \Obj(C)$ has a \emph{identity} (or \emph{unit}) morphism $\mathrm{id}_X : X \to X$ such that for every morphism $f: X \to Y$ and $g: Z \to X$:
              \begin{align}
                  \mathrm{id}_Y \circ f \;=\; f \\
                  f \circ \mathrm{id}_X \;=\; f
              \end{align}
    \end{enumerate}
\end{defB}

This definition is very abstract, so the following examples should make the idea more concrete.


\begin{expB}{Common categories}
    Here are some common categories: (See \refdef{category})
    \begin{defb}[catset]{$\catset$, the category of sets}
        \begin{itemize}
            \item $\Obj$: All sets.
            \item $\Hom$: All functions between sets.
            \item $\circ$: Normal function composition.
            \item $\id_X$: The identity function on a set $X$.
        \end{itemize}
    \end{defb}
    \begin{defb}[catrel]{$\catrel$, the category of relations}
        \begin{itemize}
            \item $\Obj$: All sets.
            \item $\Hom$: A morphism $R : X \to Y$ is a binary relation $R \subseteq X \times Y$.
            \item $\circ$ :  For $R : X \to Y$ and $S : Y \to Z$, their composite $S \circ R : X \to Z$ is given by:
                  $$
                      (x,z) \in (S \circ R) \quad \text{iff} \quad \exists\, y \in Y \text{ such that } (x,y)\in R \text{ and } (y,z)\in S.
                  $$
            \item $\id_X$:  For each set $X$, $\mathrm{id}_X = \{(x,x) : x \in X\}$, i.e.\ the diagonal relation.
        \end{itemize}
    \end{defb}
    \begin{defb}[cattypes]{$\cattyp$, the category of types}
        \begin{itemize}
            \item $\Obj$ : Types in a given type theory or programming language.

                  (think \texttt{string}, \texttt{int}, \texttt{double}, \texttt{string[]})
            \item $\Hom$ :  Pure functions (i.e. programs with no side effects) from the input type to the output type (think $\text{length} : \text{string} \to \text{number}$). Note only pure functions are allowed, so something that has side effects would not work.
            \item $\circ$ : Composition of functions.
            \item $\id_X$ : Identity function on type $X$.
        \end{itemize}
    \end{defb}

    \begin{defb}[catgroup]{$\catgrp$, the category of groups}
        \begin{itemize}
            \item $\Obj$: All groups.
            \item $\Hom$: Group homomorphisms. Formally, a morphism between groups $f : (G,\cdot^G) \to (H, \cdot^H)$ respects the group operation: $f(g_1 \cdot^G g_2) = f(g_1)\cdot^H f(g_2)$.
            \item $\circ$ : Composition is the usual function composition of homomorphisms.
            \item $\id_{(G,\cdot)}$: The identity on a group $(G,\cdot)$ is the identity homomorphism $\mathrm{id}_G$ which sends each $g \in G$ to itself.
        \end{itemize}
    \end{defb}

    The following example shows that morphisms on a category need not reduce to relations on the class of objects:

    \begin{defb}{Category of Dimensions of Linear maps}
        \begin{itemize}
            \item $\Obj$: The natural numbers $\omega$.
            \item $\Hom$: A morphism $R : n \to m$ is a $n\times m$ matrix $\in \R^{n \times m}$.
            \item $\circ$: matrix multiplication.
            \item $\id_n$: For $n \in \omega$, $\id_n$ is simply the $n \times n$ identity matrix.
        \end{itemize}
    \end{defb}
\end{expB}

\subsection{Basic Categorical Structures }
Here are some basic categorical structures that we will use, and these should give some intuition as to what a categorical definition looks like.


\subsubsection{Span and Cospan}
Maybe the easiest objects to define.

\begin{defB}[span-cospan]{Span and cospan}
    For a category $C$:
    \begin{defb}[span]{Span}

        A span consists of a diagram of the form:
        % https://q.uiver.app/#q=WzAsMyxbMSwwLCJCIl0sWzAsMCwiQSJdLFsyLDAsIkMiXSxbMCwxLCJmIiwyXSxbMCwyLCJnIl1d
        \[\begin{tikzcd}[ampersand replacement=\&]
                A \& B \& C
                \arrow["f"', from=1-2, to=1-1]
                \arrow["g", from=1-2, to=1-3]
            \end{tikzcd}\]
    \end{defb}
    \begin{defb}[cospan]{Cospan}

        A cospan consists of a diagram of the form:
        % https://q.uiver.app/#q=WzAsMyxbMSwwLCJCIl0sWzAsMCwiQSJdLFsyLDAsIkMiXSxbMSwwLCJmIl0sWzIsMCwiZyIsMl1d
        \[\begin{tikzcd}[ampersand replacement=\&]
                A \& B \& C
                \arrow["f", from=1-1, to=1-2]
                \arrow["g"', from=1-3, to=1-2]
            \end{tikzcd}\]

    \end{defb}

\end{defB}


\subsubsection{Monoid}
If we look at a category, sometimes we can see an isolated object, i.e. something that looks like this:
% https://q.uiver.app/#q=WzAsMSxbMCwwLCJBIl0sWzAsMCwiZyJdLFswLDAsImYiLDAseyJhbmdsZSI6LTQ1fV0sWzAsMCwiayIsMCx7ImFuZ2xlIjotOTB9XV0=
\[\begin{tikzcd}
        A
        \arrow["g", from=1-1, to=1-1, loop, in=55, out=125, distance=10mm]
        \arrow["f", from=1-1, to=1-1, loop, in=100, out=170, distance=10mm]
        \arrow["k", from=1-1, to=1-1, loop, in=145, out=215, distance=10mm]
    \end{tikzcd}\]

where there are no arrows going from $A$ outside of $A$. This is called a monoid. More formally:
\begin{defB}[monoid]{Monoid}
    A monoid in a category $C$ is an object $M$ such that every morphism $f$ with domain $M$ must also have codomain $M$.
\end{defB}
Recall the traditional definition of a monoid consists of:
\begin{defb}{Set theoretic monoid}
    A monoid $(M,+,e)$ is a set $M$, function $+: M\times M \to M$, element $e \in M$ s.t.:
    \begin{enumerate}
        \item $+$ is associative
        \item For all $m \in M$:
              $e + m = m + e = m$
    \end{enumerate}
\end{defb}

\begin{resb}{Category Theory Monoid is Set Monoid}
    Let $* \in \Obj (C)$ be a monoid. Then $(\Hom (*, *) , \circ)$ is a set theoretic monoid. (recall $\circ$ is the composition operator of $C$)

    \begin{proof}
        Associativity and the existence of identity follow immediatly from the category theory axioms.
    \end{proof}
\end{resb}

\subsubsection{Isomorphism and Isomorphic objects}

An isomorphism generalizes the idea of a bijection. Namely:
\begin{defB}[isomorphism]{Isomorphism}
    For a category $C$, an isomorphism is any morphism $f: X \to Y$ such that there exists another morphism $f^{-1} : Y \to X$ such that $ f^{-1} \circ f = \id_X $, and $f \circ f^{-1} = \id_Y$. As a diagram:
    % https://q.uiver.app/#q=WzAsMixbMCwwLCJYIl0sWzEsMCwiWSJdLFswLDEsImYiLDAseyJjdXJ2ZSI6LTF9XSxbMSwwLCJmXnstMX0iLDAseyJjdXJ2ZSI6LTF9XV0=
    \[\begin{tikzcd}[ampersand replacement=\&]
            X \& Y
            \arrow["f", curve={height=-6pt}, from=1-1, to=1-2]
            \arrow["{f^{-1}}", curve={height=-6pt}, from=1-2, to=1-1]
        \end{tikzcd}\]
\end{defB}
We can transfer this definition from morphisms to objects:
\begin{defB}[isomorphic]{Isomorphic objects}
    In a category $C$, objects $X,Y \in C$ are isomorphic (denoted $X\cong Y$) if there exists some isomorphism $f: X \to Y$.
\end{defB}

\begin{expb}{Isomorphisms/Isomorphic objects}
    \begin{itemize}
        \item In $\catset$, isomorphisms are exactly bijective functions, and two sets are isomorphic iff they have the same cardinality.
        \item In $\catgrp$, isomorphims are group homomorphisms that are invertible, i.e. group isomorphisms.
    \end{itemize}
\end{expb}

Many categorical definitions are only unique up to isomorphism, meaning that any two objects that satisfy the condition must be isomorphic. For example see \refres{terminal-unique}.



\subsubsection{Initial and Terminal/Final Objects}

\begin{defB}[initial-obj]{Initial Object}
    An object $* \in \Obj (C)$ for some category $C$ is called initial if for any other object $X \in \Obj (C)$, there is a unique morphism $! : * \to X$.
\end{defB}

\begin{defB}[terminal-obj]{Terminal/Final Object}
    An object $* \in \Obj (C)$ for some category $C$ is called terminal (aka final) if for any other object $X \in \Obj (C)$, there is a unique morphism $! : X \to *$.
\end{defB}

\begin{expB}{Initial/Terminal Objects}

    \begin{expb}{Initial for $\catset$}
        Note that if we look at $\catset$, the set with a single element (usually denoted $1 = \{0\} = \{\emptyset\}$) is terminal. (For any other set, there is a unique map that just sends everything to $0$)
    \end{expb}
    \begin{expb}{Terminal for $\catset$}
        By convention, we assume that there is a unique function from any set to the emptyset. This means that the empty set (aka $0$) is the terminal object in $\catset$.
    \end{expb}
    \begin{expb}{Terminal for $\catgrp$}
        Note that the trivial group $(\{0\}, +, 0)$ is terminal in $\catgrp$. Namely for any other group $(G,+,0)$ there is only one homomorphism to the trivial group, namely the one that sends everything to the identity.
    \end{expb}
\end{expB}


\begin{resB}[terminal-unique]{Terminal Objects are Unique up to Isomorphism}
    Let $*_0, *_1$ both be terminal objects of a category $C$. Then $*_0 \cong *_1$.
    \dline
    \begin{proof}
        Because $*_0$ is terminal, there must be some (unique) $!_0: *_1 \to *_0$. Likewise because $*_1$ is terminal, there must be some (unique) $!_1: *_0 \to *_1$. So notice that $!_0 \circ !_1 : *_0 \to *_0$, so by finality of $*_0$ the identity map is the unique morphism $: *_0 \to *_0$, so $!_0 \circ !_1 = \id_{*_0}$. The same argument holds for $!_1 \circ !_0 = \id_{*_1}$. Thus by \refdef{isomorphism} $!_0$ is an isomorphism, and so $*_0 \cong *_1$.
    \end{proof}
\end{resB}
Note that the same is true for initial objects. This means that we can say \textbf{the} [initial/terminal] object instead of \textbf{an/a} [initial/terminal] object.


\subsubsection{Product and Coproduct}

\begin{defB}[product]{Product}
    Let $A, B \in \Obj (C)$ for some category $C$. A product consists of:
    \begin{itemize}
        \item An object $A\times B \in \Obj(C)$
        \item A morphism $\pi_A : A \times B \to A$
        \item A morphism $\pi_B : A \times B \to B$
    \end{itemize}
    s.t. for any other span between $A$ and $B$: $A \lfrom{f} Q \lto{g} B$ $ \left(\refdef{span}\right)$  we have that there exists a unique map $r_Q:Q \to A \times B$ s.t. TFDC:
    % https://q.uiver.app/#q=WzAsNCxbMSwxLCJRIl0sWzEsMCwiQSBcXHRpbWVzIEIiXSxbMCwxLCJBIl0sWzIsMSwiQiJdLFswLDEsIlxcZXhpc3RzIXJfUSIsMCx7InN0eWxlIjp7ImJvZHkiOnsibmFtZSI6ImRhc2hlZCJ9fX1dLFswLDIsImYiXSxbMCwzLCJnIiwyXSxbMSwyLCJcXHBpX0EiLDIseyJjdXJ2ZSI6Mn1dLFsxLDMsIlxccGlfQiIsMCx7ImN1cnZlIjotMn1dXQ==
    \[\begin{tikzcd}[ampersand replacement=\&]
            \& {A \times B} \\
            A \& Q \& B
            \arrow["{\pi_A}"', curve={height=12pt}, from=1-2, to=2-1]
            \arrow["{\pi_B}", curve={height=-12pt}, from=1-2, to=2-3]
            \arrow["{\exists!r_Q}", dashed, from=2-2, to=1-2]
            \arrow["f", from=2-2, to=2-1]
            \arrow["g"', from=2-2, to=2-3]
        \end{tikzcd}\]
\end{defB}

Intuitively, this means that if I know a product exists, and I have some span $A \lfrom{f} Q \lto{g} B$, we have that $f = \pi_A \circ r_Q$, and $g = \pi_B \circ r_Q$. This means that we can factor/reduce/simplify $f,g$ by going through the product.


\textbf{IMPORTANT NOTE}: Not all categories have a product for every two objects, however many of the ones we will see do.

\begin{expb}{Products of some categories}
    The product for objects $A,B \in \Obj(C)$ in the categories $C$ below is given by:
    \begin{expb}{$\catset$}
        The product of sets $A, B$ consists of the cartesian product $\{(a,b) | a\in A, b \in B\}$, where $\pi_A, \pi_B$ are the projection functions from $A\times B$ to $A$ and $B$ respectfully. Note that given a span $A \lfrom{f} Q \lto{g} B$, the unique $r_Q \left(\text{see }\refdef{product}\right)$ is given by $r_Q = (q: Q)\mapsto (f(q), g(q))$.
    \end{expb}
    \begin{expb}{$\catgrp$} The product of groups $(G_0, +_0,e_0)$, $(G_0, +_1,e_1)$ is a group $(G_0 \times G_1, (+_0, +_1), (e_0, e_1))$, where $G_0 \times G_1$ is the cartesian product on sets, and $(+_0, +_1)$ acts elementwise.
    \end{expb}
\end{expb}


\begin{defB}[coproduct]{Coproduct}
    Let $A, B \in \Obj (C)$ for some category $C$. A coproduct consists of:
    \begin{itemize}
        \item An object $A+ B \in \Obj(C)$
        \item A morphism $i_A : A \to A + B $
        \item A morphism $i_B : B \to A + B$
    \end{itemize}
    s.t. for any other cospan between $A$ and $B$: $A \lto{f} Q \lfrom{g} B$ $ \left(\refdef{cospan}\right)$  we have that there exists a unique map $r_Q:A + B \to Q$ s.t. TFDC:
    % https://q.uiver.app/#q=WzAsNCxbMSwxLCJRIl0sWzEsMCwiQSArIEIiXSxbMCwxLCJBIl0sWzIsMSwiQiJdLFsxLDAsIlxcZXhpc3RzIXJfUSIsMix7InN0eWxlIjp7ImJvZHkiOnsibmFtZSI6ImRhc2hlZCJ9fX1dLFsyLDAsImYiLDJdLFszLDAsImciXSxbMiwxLCJpX0EiLDAseyJjdXJ2ZSI6LTJ9XSxbMywxLCJpX0IiLDIseyJjdXJ2ZSI6Mn1dXQ==
    \[\begin{tikzcd}[ampersand replacement=\&]
            \& {A + B} \\
            A \& Q \& B
            \arrow["{\exists!r_Q}"', dashed, from=1-2, to=2-2]
            \arrow["{i_A}", curve={height=-12pt}, from=2-1, to=1-2]
            \arrow["f"', from=2-1, to=2-2]
            \arrow["{i_B}"', curve={height=12pt}, from=2-3, to=1-2]
            \arrow["g", from=2-3, to=2-2]
        \end{tikzcd}\]

\end{defB}

\begin{expb}{Coproduct}
    \begin{expb}{Set}
        The coproduct of sets $A,B$ consists of the disjoint union $A\sqcup B$, where the $i_A : A \to A \sqcup B$ is the inclusion map of $A$, and $i_B : B \to A \sqcup B$ is the inclusion map of $B$. Note that for any cospan $A \lto{f} Q \lfrom{g}B$, the unique $r_Q$ $\left(\text{see \refdef{coproduct}}\right)$ is given by $r_Q = (q : Q) \mapsto \begin{cases}
                f(q) & q \in A \\
                g(q) & q \in B
            \end{cases}$


    \end{expb}

\end{expb}








\subsection{Duality}
You may have noticed many of these definitions look the same, but are sort of opposite to each other. For instance the definition of a coproduct is the same as a product, but where all of the arrows are flipped. Likewise a cospan is a span but with all arrows are flipped. This notion is called duality, many times in category theory there is some definition that is of interest, and it turns out that usually, the dual definition (consisting of reversing all of the arrows) usually describes something interesting as well. When we have a name for the base thing, and we want a name for the dual, we just add the word co- in front. (coproduct, cospan, colimit, coalgebra, comonad, \ldots)


\subsection{Functors}

A functor is sort of like a homomorphism between categories. A functor should take an object to another object, and an arrow should take an arrow to another arrow. It should also preserve composition, as well as identity morphisms. More formally:

\begin{defB}[functor]{Functor}
    Let $C,D$ be two categories. A (covariant) \textit{functor} $F : C \to D$ consists of:
    \begin{itemize}
        \item A class function on objects:

              For each object $X \in \Obj(C)$, there must exist an object $FX \in \Obj(D)$

        \item A class function on morphisms:

              For each morphism $f: X\to Y$ in $\Hom(C)$, there must exist a morphism $Ff : FX \to FY$.
    \end{itemize}


    These assignments must satisfy the following conditions:
    \begin{itemize}
        \item \textit{Preservation of $\id$}:
              \begin{equation}
                  F(\id_X^C) = \id_{FX}^D \;\text{ for all} \; X \in \Obj(C)
              \end{equation}
        \item \textit{Preservation of $\circ$}:
              \begin{equation}
                  F(g \circ^C f) = F(g) \circ^D F(f)
              \end{equation}


    \end{itemize}

\end{defB}

Intuitivly, a functor basically just adds $F$ in front of everything. In particualar, if you have any diagram that commutes:
% https://q.uiver.app/#q=WzAsNCxbMSwxLCJYIl0sWzEsMCwiWSJdLFszLDEsIloiXSxbMCwxLCJcXGNkb3RzIl0sWzAsMiwiZyIsMl0sWzEsMiwiaCJdLFswLDEsImYiXSxbMywwXSxbMCwyLCJrIiwyLHsib2Zmc2V0IjotMywiY3VydmUiOjV9XV0=
\[\begin{tikzcd}
        & Y \\
        \cdots & X && Z
        \arrow["h", from=1-2, to=2-4]
        \arrow[from=2-1, to=2-2]
        \arrow["f", from=2-2, to=1-2]
        \arrow["g"', from=2-2, to=2-4]
        \arrow["k"', shift left=3, curve={height=30pt}, from=2-2, to=2-4]
    \end{tikzcd}\]

then after lifting the diagram (applying $F$ everywhere) it should still commute:
% https://q.uiver.app/#q=WzAsNCxbMSwxLCJGWCJdLFsxLDAsIkZZIl0sWzMsMSwiRloiXSxbMCwxLCJcXGNkb3RzIl0sWzAsMiwiRmciLDJdLFsxLDIsIkZoIl0sWzAsMSwiRmYiXSxbMywwXSxbMCwyLCJGayIsMix7ImN1cnZlIjo1fV1d
\[\begin{tikzcd}
        & FY \\
        \cdots & FX && FZ
        \arrow["Fh", from=1-2, to=2-4]
        \arrow[from=2-1, to=2-2]
        \arrow["Ff", from=2-2, to=1-2]
        \arrow["Fg"', from=2-2, to=2-4]
        \arrow["Fk"', curve={height=30pt}, from=2-2, to=2-4]
    \end{tikzcd}\]


\begin{defB}[endofunctor]{Endofunctor}
    An endofunctor is simply a functor that has the same codomain, namely an endofunctor on a category $C$ is a functor $F: C \to C$.
\end{defB}


\begin{expB}{Functor Examples}
    \begin{expb}{On $\catset$}, these are functors, which take a set $X$ and a morphism $f : X \to Y$ to:
        \begin{itemize}
            \item $FX = X$,  and $F(f) = f$. This is called the identity functor.
            \item $FX = X^2$, and $F(f) =(f, f) : X^2 \to Y^2$.
            \item $FX = X + A$ for some fixed set $A$, where a morphism $f : X \to Y$ maps to:
                  \begin{align}
                      F(f) & : X + A \to Y + A                                                              \\
                      F(f) & : (u : X + A) \mapsto \begin{cases} f(u) & u \in X\\ u & u \in A   \end{cases}
                  \end{align}

            \item $FX = 2^X$, and $F(f) = f[\cdot]$. (also called the covariant powerset functor)

            \item The distribution functor:


                  \begin{defB}[distribution-functor]{Distribution Functor}
                      The distribution functor on $\catset$:
                      \begin{itemize}
                          \item On objects:
                                maps a set $X$ to the set
                                \begin{equation}
                                    \Delta X = \{P : X \to [0,1]_\R  \mid \sum_{x \in X} P(x) = 1, P\text{ has finite support}\}
                                \end{equation}
                          \item On morphisms: maps a function $f: X \to Y$ to the function $\Delta f : \Delta X \to \Delta Y$, taking $P \mapsto (y : Y ) \mapsto \sum_{x \in f^{-1}[\{y\}]}P(x)$.
                      \end{itemize}

                  \end{defB}
        \end{itemize}

    \end{expb}

\end{expB}



\begin{expb}{Functors for programmers}
    If you come from a programming background, the easiest way of thinking about a functor is of a generic type. Recall $\cattyp$ from \refdef{cattypes}. A functor on $\cattyp$ consists of a generic type/way to lift types (a function $:\texttt{Type}\to\texttt{Type}$), as well as a way to lift functions. The typical example is the list functor:
    $$ : \texttt{X} \mapsto \texttt{X[]}$$
    On morphisms, given a function/program $f : \texttt{X} \to \texttt{Y}$, you get the function $\texttt{map}(f) : \texttt{X[]} \to \texttt{Y[]}$, which just applys $f$ elementwise. For example take the function: $\texttt{len} : \texttt{string} \to \texttt{Nat}$. This function takes in a string, and returns it's length as a natural number. The lifted function $\texttt{map}(\texttt{len})$ will take in a string array, and output each strings length:
    \begin{align}
        \texttt{map}(\texttt{len}) & : \texttt{string[]}\to\texttt{Nat[]}                  \\
                                   & :\texttt{["abc", "a", ""]} \mapsto \texttt{[3, 1, 0]}
    \end{align}


    If you are familiar with generic types, then a functor is exactly a generic type (usually written as \texttt{T<X>}) as well as a way to lift maps to the generic type.
\end{expb}








\subsection{Coalgebra}
As it turns out, many systems can be modeled in terms of category theory, using coalgebras (which we will define shortly). In fact the standard notion of bisimulation and trace equivalence can be modeled canonically using fairly simple definitions, and because of how general the language of category theory is, it allows us to come up with canonical notions of bisimulation for many different systems. For now we have to define what a system is.

\begin{defB}[coalgebra]{Coalgebra}
    A coalgebra for an endofunctor $F : C \to C$ is a morphism
    \begin{equation}
        c_X : X \to FX
    \end{equation}


\end{defB}
Note that there are no additional constraints. Note that if we think of $F$ adding structure to $X$, then a coalgebra can model many different forms of systems:

\begin{expb}{Coalgebras on $\catset$}
    Consider $\catset$. Here are some standard functors, along with the systems that they represent:

    \begin{itemize}
        \item $ F(X) = X $

              Deterministic system (each state transitions to exactly one next state).

        \item $ F(X) = A \times X $

              Labeled deterministic system (each state emits a label in $A$ and transitions to a unique next state).

        \item $ F(X) = 1 + X $

              System that can either terminate (landing in the singleton $1$) or transition to a next state.

        \item $ F(X) = 1 + X + X^2 $ ($+$ is disjoint union/coproduct of sets)

              System where each state can either terminate, transition deterministically, or offer two next states, equivalently a potentially infinite tree.

        \item $ F(X) = X^A $

              System where each action $a \in A$ leads to a next state (an “action-indexed” deterministic system).


        \item $F(X) = 2^X$ (also called powerset, $\P A$)


              System where each state has a set of next states, or equivalently a non deterministic transition system.



        \item $F(X) = \Delta X$ (where $\Delta$ is the distribution functor, and so $\Delta X$ is the set of distributions on $X$ which acts on morphisms via the pushforward measure)


              System with probablistic transitions.

        \item $F(X) = \R \times (\Delta X)^A$

              Markov Decision Process (MDP) with state valued reward. See here \refdef{mdp-functor} for more.
    \end{itemize}
\end{expb}


We will also need this structure:
\begin{defB}[coalgebra-hom]{Coalgebra Homomorphism}
    Let $c_X : X\to FX$, $c_Y : Y \to FY$ be two coalgebras.

    A coalgebra homomorphism is a function $\phi : X \to Y$ such that TFDC:
    % https://q.uiver.app/#q=WzAsNCxbMCwwLCJYIl0sWzEsMCwiWSJdLFswLDEsIkZYIl0sWzEsMSwiRlkiXSxbMCwxLCJcXHBoaSJdLFsyLDMsIkZcXHBoaSIsMl0sWzAsMiwiY19YIiwyXSxbMSwzLCJjX1kiXV0=
    \[\begin{tikzcd}[ampersand replacement=\&]
            X \& Y \\
            FX \& FY
            \arrow["\phi", from=1-1, to=1-2]
            \arrow["{c_X}"', from=1-1, to=2-1]
            \arrow["{c_Y}", from=1-2, to=2-2]
            \arrow["{F\phi}"', from=2-1, to=2-2]
        \end{tikzcd}\]
\end{defB}


\begin{resB}{Coalgebras form a category}
    For a given functor $F$, the set of $F$-coalgebras form a category. Namely this category has:
    \begin{itemize}
        \item Objects: $F$-coalgebras
        \item Morphisms: $F$-coalgebra morphisms
        \item Composition: Function composition
        \item $\id_X$: Identity functions
    \end{itemize}

    \dline
    \begin{proof}
        The only things to check is that the composition of coalgebra homomorphisms is a coalgebra homomorphism, and that there are identities (which are just the identity functions).


        Assume you have 3 coalgebras $c_X, c_Y, c_Z$. Let $\phi : c_X \to c_Y$, and $\psi : c_Y \to c_Z$.
        % https://q.uiver.app/#q=WzAsNixbMCwwLCJYIl0sWzEsMCwiWSJdLFswLDEsIkZYIl0sWzEsMSwiRlkiXSxbMiwwLCJaIl0sWzIsMSwiRloiXSxbMCwxLCJcXHBoaSJdLFsyLDMsIkZcXHBoaSIsMl0sWzAsMiwiY19YIiwyXSxbMSwzLCJjX1kiXSxbMSw0LCJcXHBzaSJdLFszLDUsIkZcXHBzaSIsMl0sWzQsNSwiY19aIl1d
        \[\begin{tikzcd}[ampersand replacement=\&]
                X \& Y \& Z \\
                FX \& FY \& FZ
                \arrow["\phi", from=1-1, to=1-2]
                \arrow["{c_X}"', from=1-1, to=2-1]
                \arrow["\psi", from=1-2, to=1-3]
                \arrow["{c_Y}", from=1-2, to=2-2]
                \arrow["{c_Z}", from=1-3, to=2-3]
                \arrow["{F\phi}"', from=2-1, to=2-2]
                \arrow["{F\psi}"', from=2-2, to=2-3]
            \end{tikzcd}\]
        The diagram commutes because both $\phi$ and $\psi$ are coalgebras.

        Now because $F$ is a functor, $F\phi \circ F\psi = F(\phi \circ \psi)$, and so you can make the homomorphism : $c_X \to c_Z$:
        % https://q.uiver.app/#q=WzAsNCxbMCwwLCJYIl0sWzIsMCwiWiJdLFswLDEsIkZYIl0sWzIsMSwiRloiXSxbMCwxLCJcXHBzaSBcXGNpcmNcXHBoaSJdLFsyLDMsIkYoXFxwc2kgXFxjaXJjIFxccGhpKSIsMl0sWzAsMiwiY19YIiwyXSxbMSwzLCJjX1oiXV0=
        \[\begin{tikzcd}[ampersand replacement=\&]
                X \&\& Z \\
                FX \&\& FZ
                \arrow["{\psi \circ\phi}", from=1-1, to=1-3]
                \arrow["{c_X}"', from=1-1, to=2-1]
                \arrow["{c_Z}", from=1-3, to=2-3]
                \arrow["{F(\psi \circ \phi)}"', from=2-1, to=2-3]
            \end{tikzcd}\]

        The identity part can be shown as well by just assuming $\phi = \id^X, Y=X$ or $\psi = \id^Y, Z =Y$. Associativity can be shown as well because coalgebra homomorphism composition is the same as function composition.

    \end{proof}
\end{resB}



\newpage
\section{Slice and product category}
Later we will use the following constructions:

\begin{defB}[slice-cat]{Slice Category}
    Let $C$ be some category, and let $A \in \Obj (C)$. The slice category $C/A$ is the category with:
    \begin{itemize}
        \item $\Obj(C/A)$:
              $V_X$, where $V_X : X \to A$ is a morphism in $C$.

        \item $\Hom(C/A)$: $f : V_X \to V_Y$ reduces to $f : X\to Y$ s.t. TFDC:
              % https://q.uiver.app/#q=WzAsMyxbMCwwLCJYIl0sWzIsMCwiWSJdLFsxLDEsIkEiXSxbMCwxLCJmIl0sWzAsMiwiVl9YIiwyXSxbMSwyLCJWX1kiXV0=
              \[\begin{tikzcd}[ampersand replacement=\&]
                      X \&\& Y \\
                      \& A
                      \arrow["f", from=1-1, to=1-3]
                      \arrow["{V_X}"', from=1-1, to=2-2]
                      \arrow["{V_Y}", from=1-3, to=2-2]
                  \end{tikzcd}\]





    \end{itemize}

\end{defB}

\begin{expb}{$\catset / 2$}
    $\catset / 2$ is the category with:
    \begin{itemize}
        \item $\Obj(\catset / 2) \eqdef \{X \mid X \in \Obj(\catset)\}$
        \item $\Hom(\catset / 2) \eqdef \{m_X : X \to 2 \mid X \in \Obj(\catset)\}$
        \item $\id_X, \circ$ just like in $\catset$.
    \end{itemize}


    Note that objects in this category are just binary partitions of $\Obj(\catset)$, and morphisms are just functions that preserve the partition.

    In particular if you have an object $m_X : X \to 2$, this induces a partition $X = m_X^{-1}[\{0\}] + m_X^{-1}[\{1\}]$.

    Furthermore any morphism $f: m_X \to m_Y$ must preserve the partition, i.e. $m_Y(f(x)) = m_X(x)$ for all $x \in X$.
\end{expb}

\begin{defB}[product-cat]{Product Category}
    Let $C$ and $D$ be two categories. Define the product category $C \times D$ as the category with:

    \begin{itemize}
        \item $\Obj(C \times D) \eqdef \{(X,Y) \mid X \in \Obj(C), Y \in \Obj(D)\}$ (as a class, not necessarily a set)
        \item $\Hom (C \times D) = \{(f,g) | f \in \Hom (C), g \in \Hom(D)\}$
    \end{itemize}
\end{defB}


\begin{expb}{$\catset \times \catset$}
    $\catset \times \catset$ is the category with:
    \begin{itemize}
        \item $\Obj(\catset \times \catset) \eqdef $ all $(X,Y)$ where $X, Y \in \Obj(\catset)$.
        \item $\Hom(\catset \times \catset) \eqdef $ all $(f,g)$ where $f, g \in \Hom(\catset)$.
        \item $\circ$:  $(f,g) \circ (h,k) \eqdef (f \circ h, g \circ k)$
        \item $\id_X$: $\id_{(X,Y)} \eqdef(\id_X, \id_Y)$
    \end{itemize}

\end{expb}





\newpage
\section{Final Coalgebras, Behavioral Equivalence, Bisimulation}

\subsection{Final Coalgebras}
One natural question to answer is given two systems, are they behaviorally equivalent. This question can be formalized by looking at terminal objects in the category of coalgebras.

\begin{defb}{Terminal/Final Coalgebra}
    Recall \refdef{terminal-obj}.

    Fix a functor $F$.
    The terminal/final coalgebra for $F$ is the terminal object in the category of $F$-coalgebras.
\end{defb}

\begin{expb}{Terminal Coalgebra}
    As it turns out, most categories have a notion of a terminal object. Here we give some examples:
    %https://link.springer.com/chapter/10.1007/BFb0018361
    \begin{expb}{$FX = A \times X$}
        The final coalgebra consists of $* : A^\omega \to A \times A^\omega$, which maps $(a_0, a_1, \ldots) \mapsto (a_0, (a_1, \ldots))$.

        Any coalgebra $c: X \to A \times X$  has a unique coalgebra morphism $! : c \to *$, namely the map which sends $x$ to the sequence $(a_0, a_1, \ldots ) : A^\omega$ that happens after applying $c$ infinitely often:

        $c(x) = (a_0, x_1), c(x_1) = (a_1, x_2), \ldots$.
    \end{expb}

    \begin{expb}{$FX = 2 \times X^A$}
        The final coalgebra consists of $2^{A^*}$, the set of languages over $A$.
    \end{expb}



\end{expb}


\begin{resB}{Lambek's theorem (dual)}
    If $F$ has a final coalgebra $\alpha : A \to FA$, then $A \cong FA$.

    \dline

    \begin{proof}
        By finality of $\alpha$, there must be a unique coalgebra hom $f: F \alpha \to \alpha$. Thus TFDC:
        % https://q.uiver.app/#q=WzAsNCxbMCwwLCJBIl0sWzAsMSwiRkEiXSxbMSwxLCJGRkEiXSxbMSwwLCJGQSJdLFswLDEsIlxcYWxwaGEiLDJdLFszLDIsIkZcXGFscGhhIl0sWzIsMSwiRmYiLDAseyJzdHlsZSI6eyJib2R5Ijp7Im5hbWUiOiJkYXNoZWQifX19XSxbMywwLCJcXGV4aXN0cyAhIGYiLDIseyJzdHlsZSI6eyJib2R5Ijp7Im5hbWUiOiJkYXNoZWQifX19XV0=
        \[\begin{tikzcd}[ampersand replacement=\&]
                A \& FA \\
                FA \& FFA
                \arrow["\alpha"', from=1-1, to=2-1]
                \arrow["{\exists ! f}"', dashed, from=1-2, to=1-1]
                \arrow["{F\alpha}", from=1-2, to=2-2]
                \arrow["Ff", dashed, from=2-2, to=2-1]
            \end{tikzcd}\]
        Now note that $\alpha$ is also a coalgebra hom $:\alpha \to F\alpha$ (trivially TFDC):
        % https://q.uiver.app/#q=WzAsNCxbMCwwLCJBIl0sWzAsMSwiRkEiXSxbMSwxLCJGRkEiXSxbMSwwLCJGQSJdLFswLDEsIlxcYWxwaGEiLDJdLFswLDMsIlxcYWxwaGEiXSxbMSwyLCJGXFxhbHBoYSIsMl0sWzMsMiwiRlxcYWxwaGEiXV0=
        \[\begin{tikzcd}[ampersand replacement=\&]
                A \& FA \\
                FA \& FFA
                \arrow["\alpha", from=1-1, to=1-2]
                \arrow["\alpha"', from=1-1, to=2-1]
                \arrow["{F\alpha}", from=1-2, to=2-2]
                \arrow["{F\alpha}"', from=2-1, to=2-2]
            \end{tikzcd}\]
        So there are two coalgebra homomorphisms $:\alpha \to \alpha$, namely $\id_A$ and $f \circ \alpha$. By finality of $\alpha$, we have $f \circ \alpha = \id_A$.

        So because $f\circ \alpha = \id_A$, it must be that $Ff \circ F\alpha = \id_{FA}$, by functorality of $F$.

        Thus by the first diagram, we have $\id_{FA} = Ff \circ F\alpha = \alpha \circ f$.
    \end{proof}

\end{resB}
Note that Lambek's Theorem is for $F$ algebras, but the dual statement is what we proved above.



\subsection{Behavioral Equivalence and Bisimulation}

There are two notions of coalgebraic bisimulation that we cover:

\begin{defB}[beq-bisimulation]{Behavioral Equivalence Bisimilarity ($\catset$)}
    Define the bisimilarity relation between two $F$-coalgebras:
    $$c_X : X \to FX$$
    $$c_Y : Y \to FY$$

    as a relation on $X \times Y$ given by:
    \begin{align}
        x \sim y \iff [x]_{\texttt{final}(c_X)} = [y]_{\texttt{final}(c_Y)}
    \end{align}
    where $[\cdot]_{\texttt{final}(c_X)}$ and $[\cdot]_{\texttt{final}(c_Y)}$ are the unique homomorphisms to the final $F$-coalgebra from $c_X$ and $c_Y$.


\end{defB}

\begin{resb}{Behavioral Equivalence on a set is an equivalence relation}
    If $X = Y$ and $c_X = c_Y$, then this will be an equivalence relation.
\end{resb}

\begin{defB}[AM-bisimulation]{Aczel-Mendler/Span Bisimulation}
    An Azcel Mendler bisimulation (AM-bismulation) is a \refdef{span} in the category of $F$-coalgebras.

    I.e. for two coalgebras $c_X : X\to FX$, and $c_Y : Y \to FY$, an AM bisimulation is a span:
    % https://q.uiver.app/#q=WzAsMyxbMSwwLCJjX1IiXSxbMCwwLCJjX1giXSxbMiwwLCJjX1kiXSxbMCwxLCJcXHBoaV9YIiwyXSxbMCwyLCJcXHBoaV9ZIl1d
    \[\begin{tikzcd}[ampersand replacement=\&]
            {c_X} \& {c_R} \& {c_Y}
            \arrow["{\phi_X}"', from=1-2, to=1-1]
            \arrow["{\phi_Y}", from=1-2, to=1-3]
        \end{tikzcd}\]

\end{defB}
Note that the two are related. In fact for most common functors (any functor that preserves weak pullbacks) AM Bisimulation is weaker than Behavioral equivalence.

The result is a bit more complicated, but can be found here \refres{AM-vs-beh-equiv}.




\newpage
\section{Couplings (new)}
\begin{defb}[set-element-coupling]{Set Element Coupling}

    Let
    \begin{itemize}
        \item $F : \catset \to \catset$,
        \item $n \in \omega$,
        \item $X \in \Obj (\catset)$,
        \item $t_i \in FX$ for all $i \in n$.
    \end{itemize}
    Then $t \in F(X^n)$ is a \textit{set element coupling} iff for all $i$ we have $F\pi_i t = t_i$. Or equivalently if TFDC:
    % https://q.uiver.app/#q=WzAsMyxbMCwwLCIqIl0sWzEsMCwiRlhebiJdLFsxLDEsIkZYIl0sWzEsMiwiRlxccGlfaSJdLFswLDIsInRfaSIsMl0sWzAsMSwidCJdXQ==
    \[\begin{tikzcd}[ampersand replacement=\&]
            {*} \& {FX^n} \\
            \& FX
            \arrow["t", from=1-1, to=1-2]
            \arrow["{t_i}"', from=1-1, to=2-2]
            \arrow["{F\pi_i}", from=1-2, to=2-2]
        \end{tikzcd}\]

\end{defb}

\subsubsection{New definitions of couplings (new)}
We can generalize this definition so that it does not need to operate on set, also so that it acts on objects instead of elements, and not limit it to one object $X$ but instead a collection of different objects:

\begin{defB}[coupling]{Coupling}
    A \textit{coupling} for $\prod_{i \in I} F(X_i)$ is a morphism $c: \prod_{i \in I} F(X_i) \to F(\prod_{i \in I} X_i)$ such that TFDC:

    % https://q.uiver.app/#q=WzAsNSxbMCwzLCJcXHByb2Rfe2kgXFxpbiBJfUYoWF9pKSJdLFs0LDMsIkZcXGxlZnQoXFxwcm9kX3tpIFxcaW4gSX1YX2lcXHJpZ2h0KSJdLFsyLDIsIkZfe1hfe2lfMH19Il0sWzIsMSwiRl97WF97aV8xfX0iXSxbMiwwLCJcXGNkb3RzIl0sWzAsMiwiXFxwaV97aV8wfSIsMV0sWzEsMiwiRihcXHBpX3tpXzB9KSIsMV0sWzAsMywiXFxwaV97aV8xfSIsMSx7ImxhYmVsX3Bvc2l0aW9uIjo0MCwiY3VydmUiOi0xfV0sWzEsMywiRihcXHBpX3tpXzF9KSIsMSx7ImN1cnZlIjoxfV0sWzEsNCwiIiwwLHsiY3VydmUiOjN9XSxbMCw0LCIiLDAseyJjdXJ2ZSI6LTN9XSxbMCwxLCJjIiwxXV0=
    \[\begin{tikzcd}[ampersand replacement=\&]
            \&\& \cdots \\
            \&\& {F_{X_{i_1}}} \\
            \&\& {F_{X_{i_0}}} \\
            {\prod_{i \in I}F(X_i)} \&\&\&\& {F\left(\prod_{i \in I}X_i\right)}
            \arrow[curve={height=-18pt}, from=4-1, to=1-3]
            \arrow["{\pi_{i_1}}"{description, pos=0.4}, curve={height=-6pt}, from=4-1, to=2-3]
            \arrow["{\pi_{i_0}}"{description}, from=4-1, to=3-3]
            \arrow["c"{description}, from=4-1, to=4-5]
            \arrow[curve={height=18pt}, from=4-5, to=1-3]
            \arrow["{F(\pi_{i_1})}"{description}, curve={height=6pt}, from=4-5, to=2-3]
            \arrow["{F(\pi_{i_0})}"{description}, from=4-5, to=3-3]
        \end{tikzcd}\]

    i.e. $\forall_i^I \pi_i = F\pi_i \circ c$
\end{defB}

this relates to the earlier definition above:
\begin{resb}{Coupling $\implies$ Coupling (original) for every object}
    Let $c$ be a coupling \refdef{coupling} on $\catset$  between $X_i = X$ for all $i \in I$ for a finite $I$ (so $|I|\in \omega$). Then for any $(t_i)_{i \in I} \in \prod_{i \in I} F(X)$ is a set element coupling \refdef{set-element-coupling}.

    \dline
    \begin{proof}
        Take $p \eqdef (t_i)_{i \in I} \in \prod_{i \in I} F(X)$. From the definition of coupling, we have that $\pi_i(p) = (F(\pi_i) \circ c)p = (F(\pi_i) (c(p))$. Letting $t \eqdef c(p)$, we get $F(\pi_i) t = F\pi_i t$. Letting $n  = |I| \in \omega$ we get exactly the definition of a set element coupling.
    \end{proof}

\end{resb}

(note this is a simpler version of a distributive law for a monad, however we do not assume $F$ is a monad, and also specify that the other functor is the product functor)

Suppose now that $c_R : R \to FR$ is a bisimulation between $c_X : X \to FX$, and $c_Y : Y \to FY$. Then note because of \refdef{coalgebra-hom}, there must be some span on sets $X \lfrom{\phi_X} R \lto{\phi_Y} Y$, and because $\catset$ has products, by \refdef{product} these must factor out into product maps. The question is, is there any way to given a bisimulation $c_R : R \to FR$, reduce it to a coalgebra of type $X \times Y \to F(X \times Y)$? this would allow for an easier way to interpret what the bisimulation is doing.



\begin{resB}{Products of Coalgberas vs Products on Original Category}
    Let $C$ be a category with products.

    If a coupling $\gamma: FX \times FY \to F(X \times Y)$ exists, then for any two coalgebras $c_x : X \to FX, c_Y : Y \to FY$ we have:
    \begin{itemize}
        \item The product of $c_X, c_Y$ exists in the category of $F$-coalgebras.
        \item This product can be of type $c_{X \times Y} : X \times Y \to F(X \times Y)$
    \end{itemize}
    \dline
    \begin{proof}
        Let $c_X : X \to FX$, and $c_Y : Y \to FY$, and let $c_R: R \to FR$ be a bisimulation, with homomorphisms $\phi_X : R \to X$, $\phi_Y : R \to Y$.
        The goal is to show that a product exists, i.e. $\exists c_{X \times Y} : X \times Y \to F(X \times Y)$ such that for any bisimulation $c_R$, TFDC (in the category of F-coalgebras):

        % https://q.uiver.app/#q=WzAsNCxbMiwyLCJjX1IiXSxbMCwyLCJjX1giXSxbNCwyLCJjX1kiXSxbMiwwLCJjX3tYIFxcdGltZXMgWX0iXSxbMCwxLCJcXHBoaV9YIl0sWzAsMiwiXFxwaGlfWSIsMl0sWzMsMSwiXFxwaV9YIiwyXSxbMywyLCJcXHBpX1kiXSxbMCwzLCJcXGV4aXN0cyAhIiwxLHsic3R5bGUiOnsiYm9keSI6eyJuYW1lIjoiZGFzaGVkIn19fV1d
        \[\begin{tikzcd}[ampersand replacement=\&]
                \&\& {c_{X \times Y}} \\
                \\
                {c_X} \&\& {c_R} \&\& {c_Y}
                \arrow["{\pi_X}"', from=1-3, to=3-1]
                \arrow["{\pi_Y}", from=1-3, to=3-5]
                \arrow["{\exists !}"{description}, dashed, from=3-3, to=1-3]
                \arrow["{\phi_X}", from=3-3, to=3-1]
                \arrow["{\phi_Y}"', from=3-3, to=3-5]
            \end{tikzcd}\]

        Now going back to $C$, notice that since $C$ has products, there exists an $f : R \to X \times Y$ such that TFDC (in $C$):
        % https://q.uiver.app/#q=WzAsNCxbMiwyLCJSIl0sWzAsMiwiWCJdLFs0LDIsIlkiXSxbMiwwLCJYIFxcdGltZXMgWSJdLFswLDEsIlxccGhpX1giXSxbMCwyLCJcXHBoaV9ZIiwyXSxbMywxLCJcXHBpX1giLDJdLFszLDIsIlxccGlfWSJdLFswLDMsIlxcZXhpc3RzICFmIiwxLHsic3R5bGUiOnsiYm9keSI6eyJuYW1lIjoiZGFzaGVkIn19fV1d
        \[\begin{tikzcd}[ampersand replacement=\&]
                \&\& {X \times Y} \\
                \\
                X \&\& R \&\& Y
                \arrow["{\pi_X}"', from=1-3, to=3-1]
                \arrow["{\pi_Y}", from=1-3, to=3-5]
                \arrow["{\exists !f}"{description}, dashed, from=3-3, to=1-3]
                \arrow["{\phi_X}", from=3-3, to=3-1]
                \arrow["{\phi_Y}"', from=3-3, to=3-5]
            \end{tikzcd}\]
    \end{proof}




    Now note that we assume $c_R$ is a bisimulaiton, so we have TFDC (in $C$)

    % https://q.uiver.app/#q=WzAsNixbMSwwLCJSIl0sWzAsMCwiWCJdLFsyLDAsIlkiXSxbMiwxLCJGWSJdLFswLDEsIkZYIl0sWzEsMSwiRlIiXSxbMCwxLCJcXHBoaV9YIiwyXSxbMCwyLCJcXHBoaV9ZIl0sWzAsNSwiY19SIiwyXSxbNSw0LCJGXFxwaGlfWCJdLFs1LDMsIkZcXHBoaV9ZIiwyXSxbMSw0LCJjX1kiLDJdLFsyLDMsImNfWSJdXQ==
    \[\begin{tikzcd}[ampersand replacement=\&]
            X \& R \& Y \\
            FX \& FR \& FY
            \arrow["{c_Y}"', from=1-1, to=2-1]
            \arrow["{\phi_X}"', from=1-2, to=1-1]
            \arrow["{\phi_Y}", from=1-2, to=1-3]
            \arrow["{c_R}"', from=1-2, to=2-2]
            \arrow["{c_Y}", from=1-3, to=2-3]
            \arrow["{F\phi_X}", from=2-2, to=2-1]
            \arrow["{F\phi_Y}"', from=2-2, to=2-3]
        \end{tikzcd}\]

    Notice, that if we then consider the object $FX \times FY$, this is a product object in $C$. So it must have projection maps $\pi_{FX}, \pi_{FY}$. Furthermore, since $X \times Y \lto{\pi_X} X \lto{c_X} FX$ and $X \times Y \lto{\pi_Y} Y \lto{c_Y} FY$, this forms a span of $FX$ and $FY$, so there must be a unique map $\exists ! : X \times Y \to FX \times FY$ such that TFDC:
    % https://q.uiver.app/#q=WzAsOCxbMSwxLCJSIl0sWzEsMywiRlIiXSxbMiwyLCJZIl0sWzIsNCwiRlkiXSxbMCwwLCJYIl0sWzAsMiwiRlgiXSxbNSwxLCJYXFx0aW1lcyBZIl0sWzUsMywiRlggXFx0aW1lcyBGWSJdLFswLDIsIlxccGhpX1kiLDFdLFswLDQsIlxccGhpX1giLDFdLFsxLDUsIkZcXHBoaV9YIiwxXSxbMSwzLCJGXFxwaGlfWSIsMV0sWzAsMSwiY19SIiwxLHsibGFiZWxfcG9zaXRpb24iOjMwfV0sWzQsNSwiY19YIiwxLHsibGFiZWxfcG9zaXRpb24iOjMwfV0sWzIsMywiY19ZIiwxLHsibGFiZWxfcG9zaXRpb24iOjMwfV0sWzYsNCwiXFxwaV9YIiwxLHsibGFiZWxfcG9zaXRpb24iOjMwfV0sWzYsMiwiXFxwaV9ZIiwxLHsibGFiZWxfcG9zaXRpb24iOjYwfV0sWzAsNiwiXFxleGlzdHMgISBmIiwxLHsibGFiZWxfcG9zaXRpb24iOjYwLCJzdHlsZSI6eyJib2R5Ijp7Im5hbWUiOiJkYXNoZWQifX19XSxbNyw1LCJcXHBpX3tGWH0iLDEseyJsYWJlbF9wb3NpdGlvbiI6MzB9XSxbNywzLCJcXHBpX3tGWX0iLDFdLFs2LDcsIlxcZXhpc3RzICEgPSBcXGxhbmdsZSBjX3gsIGNfeVxccmFuZ2xlIiwwLHsic3R5bGUiOnsiYm9keSI6eyJuYW1lIjoiZGFzaGVkIn19fV1d
    \[\begin{tikzcd}[ampersand replacement=\&]
            X \\
            \& R \&\&\&\& {X\times Y} \\
            FX \&\& Y \\
            \& FR \&\&\&\& {FX \times FY} \\
            \&\& FY
            \arrow["{c_X}"{description, pos=0.3}, from=1-1, to=3-1]
            \arrow["{\phi_X}"{description}, from=2-2, to=1-1]
            \arrow["{\exists ! f}"{description, pos=0.6}, dashed, from=2-2, to=2-6]
            \arrow["{\phi_Y}"{description}, from=2-2, to=3-3]
            \arrow["{c_R}"{description, pos=0.3}, from=2-2, to=4-2]
            \arrow["{\pi_X}"{description, pos=0.3}, from=2-6, to=1-1]
            \arrow["{\pi_Y}"{description, pos=0.6}, from=2-6, to=3-3]
            \arrow["{\exists ! = \langle c_x, c_y\rangle}", dashed, from=2-6, to=4-6]
            \arrow["{c_Y}"{description, pos=0.3}, from=3-3, to=5-3]
            \arrow["{F\phi_X}"{description}, from=4-2, to=3-1]
            \arrow["{F\phi_Y}"{description}, from=4-2, to=5-3]
            \arrow["{\pi_{FX}}"{description, pos=0.3}, from=4-6, to=3-1]
            \arrow["{\pi_{FY}}"{description}, from=4-6, to=5-3]
        \end{tikzcd}\]

    Now because $F$ is a functor, the following diagram must commute as well (it is just applying the functor everywhere on the product diagram for $X\times Y$, and functors preserve commutative diagrams)

    % https://q.uiver.app/#q=WzAsNCxbMiwyLCJGUiJdLFswLDIsIkZYIl0sWzQsMiwiRlkiXSxbMiwwLCJGKFggXFx0aW1lcyBZKSJdLFswLDEsIkYoXFxwaGlfWCkiXSxbMCwyLCJGKFxccGhpX1kpIiwyXSxbMywxLCJGKFxccGlfWCkiLDJdLFszLDIsIkYoXFxwaV9ZKSJdLFswLDMsIkYoZikiLDFdXQ==
    \[\begin{tikzcd}[ampersand replacement=\&]
            \&\& {F(X \times Y)} \\
            \\
            FX \&\& FR \&\& FY
            \arrow["{F(\pi_X)}"', from=1-3, to=3-1]
            \arrow["{F(\pi_Y)}", from=1-3, to=3-5]
            \arrow["{F(f)}"{description}, from=3-3, to=1-3]
            \arrow["{F(\phi_X)}", from=3-3, to=3-1]
            \arrow["{F(\phi_Y)}"', from=3-3, to=3-5]
        \end{tikzcd}\]


    So we add it to our diagram:
    % https://q.uiver.app/#q=WzAsOSxbMSwxLCJSIl0sWzEsMywiRlIiXSxbMiwyLCJZIl0sWzIsNCwiRlkiXSxbMCwwLCJYIl0sWzAsMiwiRlgiXSxbNSwxLCJYXFx0aW1lcyBZIl0sWzUsMywiRlggXFx0aW1lcyBGWSJdLFsxLDYsIkYoWFxcdGltZXMgWSkiXSxbMCwyLCJcXHBoaV9ZIiwxXSxbMCw0LCJcXHBoaV9YIiwxXSxbMSw1LCJGXFxwaGlfWCIsMV0sWzEsMywiRlxccGhpX1kiLDFdLFswLDEsImNfUiIsMSx7ImxhYmVsX3Bvc2l0aW9uIjozMH1dLFs0LDUsImNfWCIsMSx7ImxhYmVsX3Bvc2l0aW9uIjozMH1dLFsyLDMsImNfWSIsMSx7ImxhYmVsX3Bvc2l0aW9uIjozMH1dLFs2LDQsIlxccGlfWCIsMSx7ImxhYmVsX3Bvc2l0aW9uIjozMH1dLFs2LDIsIlxccGlfWSIsMSx7ImxhYmVsX3Bvc2l0aW9uIjo2MH1dLFswLDYsIlxcZXhpc3RzICEgZiIsMSx7ImxhYmVsX3Bvc2l0aW9uIjo2MCwic3R5bGUiOnsiYm9keSI6eyJuYW1lIjoiZGFzaGVkIn19fV0sWzcsNSwiXFxwaV97Rlh9IiwxLHsibGFiZWxfcG9zaXRpb24iOjMwfV0sWzcsMywiXFxwaV97Rll9IiwxXSxbNiw3LCJcXGV4aXN0cyAhID0gXFxsYW5nbGUgY194LCBjX3lcXHJhbmdsZSIsMCx7InN0eWxlIjp7ImJvZHkiOnsibmFtZSI6ImRhc2hlZCJ9fX1dLFs4LDMsIkZcXHBpX1kiLDFdLFs4LDUsIkZcXHBpX1giLDFdLFsxLDgsIkZmIiwxXV0=
    \[\begin{tikzcd}[ampersand replacement=\&]
            X \\
            \& R \&\&\&\& {X\times Y} \\
            FX \&\& Y \\
            \& FR \&\&\&\& {FX \times FY} \\
            \&\& FY \\
            \\
            \& {F(X\times Y)}
            \arrow["{c_X}"{description, pos=0.3}, from=1-1, to=3-1]
            \arrow["{\phi_X}"{description}, from=2-2, to=1-1]
            \arrow["{\exists ! f}"{description, pos=0.6}, dashed, from=2-2, to=2-6]
            \arrow["{\phi_Y}"{description}, from=2-2, to=3-3]
            \arrow["{c_R}"{description, pos=0.3}, from=2-2, to=4-2]
            \arrow["{\pi_X}"{description, pos=0.3}, from=2-6, to=1-1]
            \arrow["{\pi_Y}"{description, pos=0.6}, from=2-6, to=3-3]
            \arrow["{\exists ! = \langle c_x, c_y\rangle}", dashed, from=2-6, to=4-6]
            \arrow["{c_Y}"{description, pos=0.3}, from=3-3, to=5-3]
            \arrow["{F\phi_X}"{description}, from=4-2, to=3-1]
            \arrow["{F\phi_Y}"{description}, from=4-2, to=5-3]
            \arrow["Ff"{description}, from=4-2, to=7-2]
            \arrow["{\pi_{FX}}"{description, pos=0.3}, from=4-6, to=3-1]
            \arrow["{\pi_{FY}}"{description}, from=4-6, to=5-3]
            \arrow["{F\pi_X}"{description}, from=7-2, to=3-1]
            \arrow["{F\pi_Y}"{description}, from=7-2, to=5-3]
        \end{tikzcd}\]

    Now note that since we have a coupling $\left(\refdef{coupling}\right)$ $\gamma : FX \times FY \to F(X\times Y)$, it must commute with $\pi_X, \pi_Y, F\pi_X,F\pi_Y$ in the diagram above, so we get TFDC:

    % https://q.uiver.app/#q=WzAsOSxbMSwxLCJSIl0sWzEsMywiRlIiXSxbMiwyLCJZIl0sWzIsNCwiRlkiXSxbMCwwLCJYIl0sWzAsMiwiRlgiXSxbNSwxLCJYXFx0aW1lcyBZIl0sWzUsMywiRlggXFx0aW1lcyBGWSJdLFsxLDYsIkYoWFxcdGltZXMgWSkiXSxbMCwyLCJcXHBoaV9ZIiwxXSxbMCw0LCJcXHBoaV9YIiwxXSxbMSw1LCJGXFxwaGlfWCIsMV0sWzEsMywiRlxccGhpX1kiLDFdLFswLDEsImNfUiIsMSx7ImxhYmVsX3Bvc2l0aW9uIjozMH1dLFs0LDUsImNfWCIsMSx7ImxhYmVsX3Bvc2l0aW9uIjozMH1dLFsyLDMsImNfWSIsMSx7ImxhYmVsX3Bvc2l0aW9uIjozMH1dLFs2LDQsIlxccGlfWCIsMSx7ImxhYmVsX3Bvc2l0aW9uIjozMH1dLFs2LDIsIlxccGlfWSIsMSx7ImxhYmVsX3Bvc2l0aW9uIjo2MH1dLFswLDYsIlxcZXhpc3RzICEgZiIsMSx7ImxhYmVsX3Bvc2l0aW9uIjo2MCwic3R5bGUiOnsiYm9keSI6eyJuYW1lIjoiZGFzaGVkIn19fV0sWzcsNSwiXFxwaV97Rlh9IiwxLHsibGFiZWxfcG9zaXRpb24iOjMwfV0sWzcsMywiXFxwaV97Rll9IiwxXSxbNiw3LCJcXGV4aXN0cyAhID0gXFxsYW5nbGUgY194LCBjX3lcXHJhbmdsZSIsMCx7InN0eWxlIjp7ImJvZHkiOnsibmFtZSI6ImRhc2hlZCJ9fX1dLFs3LDgsIlxcZ2FtbWEiLDAseyJjdXJ2ZSI6LTV9XSxbOCw1LCJGXFxwaV9YIiwxXSxbMSw4LCJGZiIsMV0sWzgsMywiRlxccGlfWSIsMV1d
    \[\begin{tikzcd}[ampersand replacement=\&]
            X \\
            \& R \&\&\&\& {X\times Y} \\
            FX \&\& Y \\
            \& FR \&\&\&\& {FX \times FY} \\
            \&\& FY \\
            \\
            \& {F(X\times Y)}
            \arrow["{c_X}"{description, pos=0.3}, from=1-1, to=3-1]
            \arrow["{\phi_X}"{description}, from=2-2, to=1-1]
            \arrow["{\exists ! f}"{description, pos=0.6}, dashed, from=2-2, to=2-6]
            \arrow["{\phi_Y}"{description}, from=2-2, to=3-3]
            \arrow["{c_R}"{description, pos=0.3}, from=2-2, to=4-2]
            \arrow["{\pi_X}"{description, pos=0.3}, from=2-6, to=1-1]
            \arrow["{\pi_Y}"{description, pos=0.6}, from=2-6, to=3-3]
            \arrow["{\exists ! = \langle c_x, c_y\rangle}", dashed, from=2-6, to=4-6]
            \arrow["{c_Y}"{description, pos=0.3}, from=3-3, to=5-3]
            \arrow["{F\phi_X}"{description}, from=4-2, to=3-1]
            \arrow["{F\phi_Y}"{description}, from=4-2, to=5-3]
            \arrow["Ff"{description}, from=4-2, to=7-2]
            \arrow["{\pi_{FX}}"{description, pos=0.3}, from=4-6, to=3-1]
            \arrow["{\pi_{FY}}"{description}, from=4-6, to=5-3]
            \arrow["\gamma", curve={height=-30pt}, from=4-6, to=7-2]
            \arrow["{F\pi_X}"{description}, from=7-2, to=3-1]
            \arrow["{F\pi_Y}"{description}, from=7-2, to=5-3]
        \end{tikzcd}\]


    Thus finally we see that there is a coalgebra $c_{X \times Y} \eqdef \gamma \circ \langle c_x, c_y  \rangle : X \times Y \to F(X\times Y) $, and there are homomorphisms $f : c_R \to c_{X \times Y}$, $\pi_{X} : c_{X \times Y} \to c_X$, and $\pi_Y : c_{X \times Y} \to c_Y$.

    Now to show that the coalgebra homomorphism $f$ (not $f$ on $\catset$, but $f$ as a coalgebra homomorphism) is unique, simply note that any coalgebra homomorphism $p: c_R \to c_{X \times Y}$ must be a map $f : R \to X \times Y$, and since on $\catset$ this is unique, it must also be unique.

\end{resB}

\newpage
\section{Coalgebra of MDPs (new)}

Recall the MDP functor:
\begin{defB}[mdp-functor]{MDP functor}
    Define the MDP functor on $\catset$ as:
    \begin{equation}
        F : S \to \R \times (\Delta S)^A
    \end{equation}
\end{defB}

Note this will consider MDPs with state valued reward. Alternatively you could consider $ F : S \to (\R \times \Delta S)^A$, which assigns a reward to a transition $(s,a, s')$ instead of a state $s$. Much of the work we do should transfer up to some minor changes.


\begin{resB}[MDP-hom]{MDP homomorphim}

    Consider MDP coalgebras:
    \begin{align}
        \mathcal M_X & :X \to FX                                       \\
                     & : x \mapsto (r^X_x : \R , P_x^X : (\Delta X)^A) \\
        \mathcal M_Y & :Y \to FY                                       \\
                     & : y \mapsto (r^Y_y : \R , P_y^Y : (\Delta Y)^A)
    \end{align}

    Note that $\phi$ is a MDP-coalgebra homomorphism iff:
    % https://q.uiver.app/#q=WzAsNCxbMCwwLCJYIl0sWzEsMCwiWSJdLFsxLDEsIkZZIl0sWzAsMSwiRlgiXSxbMCwzLCJcXG1hdGhjYWwgTV9YIiwyXSxbMywyLCJGXFxwaGkiLDJdLFswLDEsIlxccGhpIl0sWzEsMiwiXFxtYXRoY2FsIE1fWSJdXQ==
    \[\begin{tikzcd}[ampersand replacement=\&]
            X \& Y \\
            FX \& FY
            \arrow["\phi", from=1-1, to=1-2]
            \arrow["{\mathcal M_X}"', from=1-1, to=2-1]
            \arrow["{\mathcal M_Y}", from=1-2, to=2-2]
            \arrow["{F\phi}"', from=2-1, to=2-2]
        \end{tikzcd}\]
    \begin{align}
        \ldots & \iff (F\phi) \circ \mathcal M_X = \mathcal M_Y \circ \phi                                                                                    \\
               & \iff \forall_x^X (F\phi) \circ \mathcal M_X (x) = \mathcal M_Y \circ \phi(x)                                                                 \\
               & \iff \forall_x^X(F\phi)(r^X_x, P_x^X) = (r_{\phi(x)}^Y, P_{\phi(x)}^Y)                                                                       \\
               & \iff \forall_x^X (r_x^X, (\Delta\phi) P_x^X) =(r_{\phi(x)}^Y, P_{\phi(x)}^Y)                                                                 \\
               & \iff \forall_x^X \left(r_x^X, (y : Y, a : A) \mapsto \sum_{\substack{x' \in X                                                                \\\phi(x') = y}}P_x^X(a,x')\right) =(r_{\phi(x)}^Y, P_{\phi(x)}^Y)\\
               & \iff \forall_x^X (r_x^X = r_{\phi(x)}^Y) \land \forall_y^Y \forall_a^A \left(\sum_{\substack{x' \in X                                        \\ \phi(x) = y}}{P_x^X (a,x')} = P_{\phi(x)}^Y (a,y) \right)\\
               & \iff \forall_x^X (r_x^X = r_{\phi(x)}^Y) \land \forall_y^Y \forall_a^A \left(P^Y(\phi(x) \lto{a} y) = P^X(x \lto{a} \phi^{-1}[\{y\}] \right)
    \end{align}

\end{resB}



\subsection{Bisimulation}
Recall the transition system version of bisimulation:
\begin{defB}[standard-mdp-bisim]{Standard MDP Bisimulation}
    For an MDP $\M = (S,A,P,R,\gamma)$, a bisimulation is an equivalence relation $B \subseteq S^2$ such that:
    \begin{align}
        s B s' \iff R(s) = R(s') \land  \forall_a^A \forall_C^{S/B} P[s \lto{a} C] = P[s' \lto{a} C]
    \end{align}
    And bisimilarity is the union of all such $B$.

\end{defB}

Now conversly, when you look at the coalgebraic version, some natural questions arise:
\begin{enumerate}
    \item How can you get a relation out of a span?
    \item If you take the relation, also assume $c_y = c_x$, does this definition correspond to the standard one?
\end{enumerate}

For question 1, the natural choice for a relation is $\{(\phi_x(b), \phi_y(b) | b \in R\}$. Or equivalently using categories, simply take the unique map $f$ to the product, and the image of that map:
% https://q.uiver.app/#q=WzAsNCxbMSwyLCJSIl0sWzAsMiwiWCJdLFsyLDIsIlkiXSxbMSwwLCJYIFxcdGltZXMgWSJdLFswLDEsIlxccGhpX1giLDJdLFswLDIsIlxccGhpX1kiXSxbMywxXSxbMywyXSxbMCwzLCJcXGV4aXN0cyAhIGYiLDIseyJzdHlsZSI6eyJib2R5Ijp7Im5hbWUiOiJkYXNoZWQifX19XV0=
\[\begin{tikzcd}
        & {X \times Y} \\
        \\
        X & B & Y
        \arrow[from=1-2, to=3-1]
        \arrow[from=1-2, to=3-3]
        \arrow["{\exists ! f}"', dashed, from=3-2, to=1-2]
        \arrow["{\phi_X}"', from=3-2, to=3-1]
        \arrow["{\phi_Y}", from=3-2, to=3-3]
    \end{tikzcd}\]

\begin{defB}[induced-relation-bisimulation]{Induced Relation for $\catset$ bisimulation}
    Let $F$ be an endofunctor on $\catset$. For a bisimulation $c_B$ of type given below:

    % https://q.uiver.app/#q=WzAsNixbMSwwLCJCIl0sWzEsMSwiRkIiXSxbMiwwLCJYIl0sWzIsMSwiRlgiXSxbMCwwLCJYIl0sWzAsMSwiRlgiXSxbMCwxLCJjX0IiXSxbMiwzLCJjX3giXSxbNCw1LCJjX1giLDJdLFswLDQsIlxccGhpXzAiLDJdLFswLDIsIlxccGhpXzEiXSxbMSwzLCJGIFxccGhpXzEiLDJdLFsxLDUsIkYgXFxwaGlfMCJdXQ==
    \[\begin{tikzcd}[ampersand replacement=\&]
            X \& B \& X \\
            FX \& FB \& FX
            \arrow["{c_X}"', from=1-1, to=2-1]
            \arrow["{\phi_0}"', from=1-2, to=1-1]
            \arrow["{\phi_1}", from=1-2, to=1-3]
            \arrow["{c_B}", from=1-2, to=2-2]
            \arrow["{c_x}", from=1-3, to=2-3]
            \arrow["{F \phi_0}", from=2-2, to=2-1]
            \arrow["{F \phi_1}"', from=2-2, to=2-3]
        \end{tikzcd}\]
    Let $B_I$ denote the induced relation on $X^2$, namely $\{\phi_X(b), \phi_X(b) | b \in B\}$.
\end{defB}

For 2, it is a bit more tricky. I think it can be done if you assume that the resulting relation is an equivalence relation.

We give the formal proof first, then give some intuition.
\begin{resB}{Coalgebraic Bisimulation is MDP bisimulation}
    Suppose $c_B$ is an MDP-coalgebraic bisimulation:
    % https://q.uiver.app/#q=WzAsMyxbMSwwLCJjX0IiXSxbMCwwLCJjX1giXSxbMiwwLCJjX1giXSxbMCwxLCJcXHBoaV97WF8wfSIsMl0sWzAsMiwiXFxwaGlfe1hfMX0iXV0=
    \[\begin{tikzcd}[ampersand replacement=\&]
            {c_X} \& {c_B} \& {c_X}
            \arrow["{\phi_{X_0}}"', from=1-2, to=1-1]
            \arrow["{\phi_{X_1}}", from=1-2, to=1-3]
        \end{tikzcd}\]
    Such that the induced relation $B_I$ is an equivalence relation. Then $B_I$ is a traditional MDP bisimulation.

    \dline
    \begin{proof}\hfill


        \begin{resb}{Forward direction}
            WTS $\forall_{x, x'} x B_I x' \implies r^X(x)  = r^X(x') \land \forall_a^A \forall_C^{X/B_I} P^X[x \lto{a} C] = P^X[x' \lto{a} C]$

            Assume $x B_I x'$.


            From \refdef{induced-relation-bisimulation}, we have $b \in B$ such that $\phi_0(b) = x$, and $\phi_1(b) = x'$.
            \begin{align}
                \implies r^B(b) & = r^X(\phi_0 (b)) = r^X(\phi_1(b)) \\
                                & = r^X(x) = r^X(x')
            \end{align}
            % rewards now equal, now show Probs are too

            Take arbitrary $a \in A$, $C \in X / B_I$. Note $\phi_0^{-1}[C] = \phi_1^{-1}[C]$, because:
            \begin{resb}{$\phi_0^{-1}[C] = \phi_1^{-1}[C]$}
                Recall $B_I = \{\phi_0(b), \phi_1(b) | b \in B\}$.

                So $\phi_0(b) \in C \iff \phi_1(b) \in C$. Thus:
                \begin{align}
                    \phi_0^{-1}[C] & = \{b \in B | \phi_0(b) \in C\} \\
                                   & = \{b \in B | \phi_1(b) \in C\} \\
                                   & = \phi_1^{-1}[C]
                \end{align}
            \end{resb}

            Using the above, and \refres{MDP-hom}:

            \begin{align}
                \forall_c^C & P^B[ b \lto{a} \phi_0^{-1}[\{c\}]] = P^X[x \lto{a} c] \\
                \implies    & P^B[ b \lto{a} \phi_0^{-1}[C]] = P^X[x \lto{a} C]     \\
                            & = P^B[ b \lto{a} \phi_1^{-1}[C]] = P^X[x' \lto{a} C]
            \end{align}
            Thus $P^X[x \lto{a} C] = P^X[x' \lto{a} C]$, and $r^X(x) = r^X(x')$.
        \end{resb}

        \begin{resb}{Backward direction} WTS $\forall_{x, x'} r^X(x)  = r^X(x') \land \forall_a^A \forall_C^{X/B_I} P^X[x' \lto{a} C] = P^X[x \lto{a} C] \implies x B_I x'$
            \todo[inline]{Prove. May not be true, but you may be able to add more conditions until you get it. This in some sense says that the bisimulation is larger,}
        \end{resb}
    \end{proof}
\end{resB}


\begin{expb}{Intuition for Bisimulation}
    Now, notice that since $B_I$ is an equivalence relation, you can view it as a matrix of size $|X| \times |X|$, where each entry corresponds to elements of $B_I \subseteq X^2$. Furthermore, you can reorder the indecies so that the matrix is block diagonal (using transitivity, reflexivity, symmetry of $B$)

    In the following figure, the red dots are elements of $B$, the green squares correspond to the actual matrix entries (so each green square corresponds to a pair $(x_0, x_1) \in B_I \subseteq X^2$, and the blue squares represent equivalence classes of $B$. Everything outside of a blue square will be empty.
    \begin{figure}[H]

        \includesvg[inkscapelatex=false, scale=0.8]{images/eq-classes.svg}
        \centering
        \caption{Bisimulation B grouped by equivalence classes}
    \end{figure}

    Notice that $P^B[b \lto{a} \phi_{X_0}^{-1} [\{x\}] = P^X[\phi_{X_0}(b)\lto{a} x]$ and $P^B[b \lto{a} \phi_{X_1}^{-1} [\{x'\}] = P^X[\phi_{X_1}(b)\lto{a} x']$ imply:
    \begin{center}
        \begin{minipage}{0.45\textwidth}
            \begin{figure}[H]

                \includesvg[inkscapelatex=false, scale=0.8]{images/eq-classes-x.svg}
                \centering
                \caption{The Transitions from any element $b\in \phi_{X_0}^{-1}[\{x\}]$, to the entirety of $\phi_{X_0}^{-1}[\{x'\}]$}
            \end{figure}
            Let $b$ be one of the red dots pictured

            (i.e. some $b\in B$ s.t. $\phi_{X_0}(b) = x$)


            Each transition above will happen with probability:
            $$P^B[b\lto{a} \phi_{X_0}[\{x'\}]] = P^X[x \lto{a} x']$$

        \end{minipage}
        \hspace{3em}
        \begin{minipage}{0.45\textwidth}

            \begin{figure}[H]

                \includesvg[inkscapelatex=false, scale=0.8]{images/eq-classes-y.svg}
                \centering
                \caption{The transitions from any element $b \in \phi_{X_1}^{-1}[\{x\}]$, to the entirety of $\phi_{X_1}^{-1}[\{x''\}]$}
            \end{figure}
            Let $b$ be one of the red dots pictured

            (i.e. some $b\in B$ s.t. $\phi_{X_1}(b) = x$)


            Each transition above will happen with probability:
            $$P^B[b\lto{a} \phi_{X_1}[\{x'\}]] = P^X[x \lto{a} x']$$

        \end{minipage}

    \end{center}
    \hspace{10em}

    Now let $C = \{q \mid q B_I x'\}$ be an equivalence class of $B_I$ that includes $x'$.

    So we can see that by summing up instances of the last picture, we get:
    \begin{figure}[H]

        \includesvg[inkscapelatex=false, scale=0.8]{images/eq-classes-b.svg}
        \centering
        \caption{Transitions from any element $b \in \phi_{X_1}^{-1}[\{x\}]$, to $\phi_{X_1}^{-1}[C]$, i.e. anything related to $x'$ }
    \end{figure}
    In the figure above, the probability of each of these transitions is $$P^B\left[b\lto{a} \phi_{X_1}^{-1}[C]\right] = P^X[x\lto{a} C]$$




    Likewise the same argument holds when you sum up columns instead, and you get
    $$P^B\left[b \lto{a} \phi_{X_0}^{-1}[C]\right] = P^X[x \lto{a} C]$$

    Note that because $B_I$ is an equivalence relation, it must be that $\phi_{X_0}^{-1}[C] = \phi_{X_1}^{-1}[C]$.

    % So any two red dots that are in the same row have the same transition probability to $\phi^{-1}_{X_1}[C]$. If they are in the same column, then they have the same probability of going to $\phi^{-1}_{X_0}[C]$. The idea is that you want to prove $x B_I x'' \implies P^X[x \lto{a}C] =P^X[x''\lto{a}C] $.
\end{expb}


\newpage
\section{LDBA/LDBW Coalgebraically (new)}

\subsection{Automota review}


Recall:
\begin{defB}[GBA]{Generalized Buchi Automoton (GBA)}
    A Generlized Buchi Automota (GBA) is a tuple

    \begin{equation}
        \B = (Q, \Sigma, \delta, \alpha)
    \end{equation}
    such that:
    \begin{itemize}
        \item $Q$ is a set (states)
        \item $\Sigma$ is a set (alphabet)
        \item $\delta : Q \to \Sigma \to 2^Q$ is a transition function (NDTM)
        \item $\alpha = \{F_1, \ldots F_n\}$ with $F_i : 2^Q$ is a set of accepting conditions
    \end{itemize}
\end{defB}

Coalgebraically, we can view a GBA as $F:  Q \to (2^Q)^\Sigma \times 2^n$,
where the left part yields the next possible states given a letter, and the
right part gives which (if any) of the acceptance conditions the current state is a part of.

\begin{defB}{Pointed GBA}
    A pointed GBA is a GBA $\B = (Q, \Sigma, \delta, \alpha)$ along with an initial state $q_0 \in Q$.
\end{defB}

For now we will consider non pointed GBA, but threre exist pointed coalgebras, so perhaps pointed GBAs reduce to pointed coalgebras.

\begin{defB}{$\inf \tau$ for trajector $\tau$}
    For a run/trajectory $\tau : (Q \times \Sigma \times Q)$, let:
    \begin{equation}
        \inf \tau \eqdef \{ q | (q, \sigma, q') \in \tau \text{ occurs infinitely often in }
        \tau\}
    \end{equation}
\end{defB}

\begin{defb}{Accepting Run}
    A run $\tau : (Q \times \Sigma \times Q)$ is accepted if $\inf \tau \circ F_i \ne 0$ for all $F_i$.
\end{defb}

\begin{defB}{Limit Deterministic Buchi Automota (LDBA)}
    A GBA is a LDBA iff $Q$ can be partitioned $Q = Q_I \sqcup Q_D$ such that:
    \begin{itemize}
        \item Tansitions from $Q_D$ must remain in $Q_D$, and must be deterministic:

              For any $q_d :  Q_D$, $\sigma : \Sigma$, $\delta(q_, \sigma) \subseteq Q_D$ and $|\delta(q, \sigma)| = 1$.
        \item Accepting states are in $Q_D$:

              For any $F_i : \alpha$, $F_i \subseteq Q_D$.
    \end{itemize}


    So the definition of a GBA can be refined for a LDBA to be:
    \begin{equation}
        (\Sigma, Q_I\sqcup Q_D, \delta_I, \delta_D, G )
    \end{equation}
    where:
    \begin{itemize}
        \item $\delta_I  : Q_I \to \Sigma \to 2^{Q_I} \times 2^{Q_D}$
        \item $\delta_D  : Q_D \to \Sigma \to Q_D$
        \item $G : Q_D \to 2^k$ (tells you which of the acceptance criteria are met)
    \end{itemize}

\end{defB}

Coalgebraically this is more tricky. Our goal is to find a functor $F$ such that every $F$-coalgebra can be modeled by an LDBA, and likewise that every LDBA can be modeled by a $F$-coalgebra.

The naiive approach would be something like: $F: Q \to (2^Q)^\Sigma + (Q^\Sigma \times 2^k)$. The problem is enforcing that transitions from $Q_D$ remain in $Q_D$.
(otherwise you can construct a coalgebra where the output of one of the deterministic transitions is something not in $Q_D$). So we need to enforce that if the coalgebra ``knows'' that the input was in $Q_D$,
it should only output elements in $Q_D$.

We can model this however by going to $\catset / 2$, i.e. \refdef{slice-cat}. This will consist of
objects of the form $m_X : X \to 2$ for any set $X$, and morphisms of the form $f: m_X \to m_Y$ which make TFDC:

% https://q.uiver.app/#q=WzAsMyxbMCwwLCJYIl0sWzIsMCwiWSJdLFsxLDEsIjIiXSxbMCwxLCJmIl0sWzAsMiwibV94IiwyXSxbMSwyLCJtX3kiXV0=
\[\begin{tikzcd}[ampersand replacement=\&]
        X \&\& Y \\
        \& 2
        \arrow["f", from=1-1, to=1-3]
        \arrow["{m_x}"', from=1-1, to=2-2]
        \arrow["{m_y}", from=1-3, to=2-2]
    \end{tikzcd}\]

Note that any object $m_X : X \to 2$ in $\catset / 2$ can be thought of as a partition $X = m_X^{-1}[\{0\}] + m_X^{-1}[\{1\}]$.

Furthermore any morphism $m_X \to m_Y$ must preserve the partition, i.e. $m_X^{-1}[\{0\}] = m_Y^{-1}[\{0\}]$ and $m_X^{-1}[\{1\}] = m_Y^{-1}[\{1\}]$.


Equivalently, we can think of $\catset / 2$ as $\catset \times \catset$. In particualr, notice that $\catset \times \catset$ is exactly $\catset / 2$, just with objects relabeled.

Namely there is a bijection:
\begin{resb}{$\catset \times \catset$ is $\catset / 2$}
    There is a bijective functor $F : \catset \times \catset \to \catset / 2$.


    Namely for any $(X,Y) : \Obj(\catset \times \catset)$, define $F$:
    \begin{itemize}
        \item On objects:
              \begin{align}
                  F(X) & : X + Y \to 2                          \\
                  F(X) & \eqdef(x : X + Y) \mapsto \begin{cases}
                                                       0 & x : X \\
                                                       1 & x : Y
                                                   \end{cases}
              \end{align}
        \item On morphisms (given $(f,g) : (X_0, X_1) \to (Y_0, Y_1)$):
              \begin{align}
                  F((f,g)) & : \Big(X_0 + X_1\Big) \to \Big(Y_0 + Y_1\Big) \\
                  F((f,g)) & \eqdef (x : X_0 + X_1) \mapsto \begin{cases}
                                                                f(x) & x : X_0 \\
                                                                g(x) & x : X_1
                                                            \end{cases}
              \end{align}

    \end{itemize}

    \dline
    \begin{proof}
        To show this is a bijection, define the inverse $G : \catset/2 \to \catset \times \catset$:
        \begin{itemize}
            \item On objects:
                  $G(m_X) \eqdef (m_X^{-1}[\{0\}], m_X^{-1}[\{1\}])  $.

            \item On morphisms:
                  $G(f : m_X \to m_Y) \eqdef ( f |_{m_{X}^{-1}[\{0\}]}, f |_{m_X^{-1}[\{1\}]})$, i.e. the restriction of $f$ to everything that maps to 0, and to 1.
        \end{itemize}

        Note that $GF = 1_{\catset \times \catset}$, and likewise $FG = 1_{\catset/2}$ for both objects and morphisms.


    \end{proof}
\end{resb}

(This is called an isomorphism of categories, it essentially means one is just relabeling the other)

This gives a more convient way to package $\catset / 2$, and also gives an easier way to define how functors act on it.

\begin{defB}{LDBA Functor}
    Consider the category $\catset \times \catset$. Define the LDBA functor such that:

    \begin{itemize}
        \item On objects:
              \begin{align}
                  F(Q_I, Q_D) & \eqdef ((2^{Q_I} \times 2^{Q_D})^\Sigma, Q_D^\Sigma \times 2^k)
              \end{align}

              Here we can see that an element of $F(X_I, X_D)$ is of the form:
              \begin{align}
                  \left[\Big((\sigma : \Sigma) \mapsto (A_\sigma : 2^{Q_I}, B_\sigma : 2^{Q_F})  \Big), \Big((\sigma : \Sigma) \mapsto (\delta : Q_D^\Sigma,G_\sigma : 2^k)\Big)
                      \right]
              \end{align}


        \item \textbf{On morphisms:}
              Given a morphism $(f : X_I \to Y_I,\; g : X_D \to Y_D)$, define
              \begin{align}
                  F(f, g) & \eqdef \left[\Big(\sigma \mapsto (A_\sigma, B_\sigma)  \Big), \Big(\sigma \mapsto (\delta,G_\sigma)\Big)
                  \right] : F(X_I, X_D)                                                                                              \\
                          & \mapsto \left[
                      \Big(\sigma \mapsto (f[A_\sigma], g[B_\sigma]) \Big), \Big(\sigma \mapsto (g \circ \delta, G_\sigma) \Big)
                      \right] : F(Y_I, Y_D)
              \end{align}
              or equivalently:

              \begin{align}
                  F(f,g) & \eqdef \left[\Big((f[\cdot], g[\cdot]) \circ (\cdot)\Big) ,   \Big(g\circ (\cdot) , \id_{2^k}\big)  \right]
              \end{align}

    \end{itemize}

\end{defB}

Notice that because we work on $\catset \times \catset$, and due to the product nature of the functor, we can decompose any LDBA coalgebra:

\begin{resb}{Equivalent LDBA coalgebra}
    Any LDBA coalgebra $c_Q : Q_I \times Q_D \to ((2^{Q_I} \times 2^{Q_D})^\Sigma, Q_D^\Sigma \times 2^k) $ must consist of two maps:

    (from definition of morphisms on $\catset \times \catset$)
    \begin{align}
        c_{Q_I} & : Q_I \to (2^{Q_I} \times 2^{Q_D})^\Sigma \\
        c_{Q_D} & : Q_D \to (Q_D^\Sigma \times 2^k)
    \end{align}

    So a LDBA coalgebra can also be defined by $(c_{Q_I}, c_{Q_D})$ with types given above.
\end{resb}

\begin{resb}{Coalgebraic LDBA model LDBA}
    There is a (class) bijection between the class of all LDBA and the objects of the category of LDBA-coalgebaras, \textbf{up to an initial state.}

    \dline
    \begin{proof} The following maps are (class) inverses of each other:

        \begin{itemize}
            \item (\textbf{LDBA $\to$ Coalgebra}) Given a LDBA (without initial state) $\B =  (\Sigma, Q_I\sqcup Q_D, \delta_I, \delta_D, G)$, you can define a LDBA coalgebra
                  with:
                  \begin{itemize}
                      \item $c_{Q_I} = \delta_I$
                      \item $c_{Q_D} =  q\mapsto(\delta_D(q),G(q))$
                  \end{itemize}

            \item (\textbf{Coalgebra $\to$ LDBA}) Given a coalgebra
                  $(c_{Q_I}, c_{Q_D}) : (Q_I , Q_D) \to F(Q_I, Q_D)$, define a LDBA as:
                  \begin{align}
                      (\Sigma, Q_I \sqcup Q_D, c_{Q_I}, \pi_0 \circ c_{Q_D}, \pi_1 \circ c_{Q_D})
                  \end{align}
        \end{itemize}
    \end{proof}

\end{resb}

\begin{resB}{LDBA Homomorphisms}
    Let:
    \begin{itemize}
        \item $c_X = (c_{X_I}, c_{X_D})$ be a LDBA coalgebra $: (X_I,X_D) \to F(X_I, X_D)$.

              Let $(X_I \sqcup X_D, \delta_{X_I}, \delta_{X_D}, G_{X})$ be the corresponding LDBA.

        \item $c_Y = (c_{Y_I}, c_{Y_D})$ be a LDBA coalgebra $: (Y_I, Y_D) \to F(Y_I , Y_D)$

              Let $(X_I \sqcup X_D, \delta_{X_I}, \delta_{X_D}, G_{X})$ be the corresponding LDBA.
    \end{itemize}
    Then a coalgbra homomorphism $\phi = (\phi_I, \phi_D) : c_X \to c_Y$ is:

    \dline
    $\phi$ will make TFDC:
    % https://q.uiver.app/#q=WzAsNCxbMCwwLCIoWF9JLFhfRCkiXSxbMiwwLCIoWV9JLFlfRCkiXSxbMCwyLCJGKFhfSSxYX0QpIl0sWzIsMiwiRihZX0ksWV9EKSJdLFswLDEsIihcXHBoaV9JLFxccGhpX0QpIl0sWzIsMywiRihcXHBoaV9JLCBcXHBoaV9EKSJdLFswLDIsIihjX3tYX0l9LCBjX3tYX0R9KSIsMl0sWzEsMywiKGNfe1lfSX0sIGNfe1lfRH0pIl1d
    \[\begin{tikzcd}[ampersand replacement=\&]
            {(X_I,X_D)} \&\& {(Y_I,Y_D)} \\
            \\
            {F(X_I,X_D)} \&\& {F(Y_I,Y_D)}
            \arrow["{(\phi_I,\phi_D)}", from=1-1, to=1-3]
            \arrow["{(c_{X_I}, c_{X_D})}"', from=1-1, to=3-1]
            \arrow["{(c_{Y_I}, c_{Y_D})}", from=1-3, to=3-3]
            \arrow["{F(\phi_I, \phi_D)}", from=3-1, to=3-3]
        \end{tikzcd}\]
    which commutes iff both of the following commute:


    \begin{minipage}{0.45\textwidth}
        % https://q.uiver.app/#q=WzAsNCxbMCwwLCJYX0kiXSxbMiwwLCJZX0kiXSxbMCwyLCIoMl57WF9JfVxcdGltZXMgMl57WF9EfSleXFxTaWdtYSJdLFsyLDIsIigyXntZX0l9XFx0aW1lcyAyXntZX0R9KV5cXFNpZ21hIl0sWzAsMSwiXFxwaGlfSSJdLFswLDIsImNfe1hfSX0iLDJdLFsxLDMsImNfe1lfSX0iXSxbMiwzLCIoXFxwaGlfSVtcXGNkb3RdLCBcXHBoaV9EW1xcY2RvdF0pIFxcY2lyYyAoXFxjZG90KSIsMl1d
        \[\begin{tikzcd}[ampersand replacement=\&]
                {X_I} \&\& {Y_I} \\
                \\
                {(2^{X_I}\times 2^{X_D})^\Sigma} \&\& {(2^{Y_I}\times 2^{Y_D})^\Sigma}
                \arrow["{\phi_I}", from=1-1, to=1-3]
                \arrow["{c_{X_I}}"', from=1-1, to=3-1]
                \arrow["{c_{Y_I}}", from=1-3, to=3-3]
                \arrow["{(\phi_I[\cdot], \phi_D[\cdot]) \circ (\cdot)}"', from=3-1, to=3-3]
            \end{tikzcd}\]
    \end{minipage}
    \begin{minipage}{0.45\textwidth}
        \vspace{1em}
        % https://q.uiver.app/#q=WzAsNCxbMCwwLCJYX0QiXSxbMiwwLCJZX0QiXSxbMCwyLCIoWF9EXlxcU2lnbWEgXFx0aW1lcyAyXmspIl0sWzIsMiwiKFlfRF5cXFNpZ21hIFxcdGltZXMgMl5rKSJdLFswLDEsIlxccGhpX0QiXSxbMCwyLCJjX3tYX0R9IiwyXSxbMSwzLCJjX3tZX0R9Il0sWzIsMywiIChcXHBoaV9EIFxcY2lyYyAoXFxjZG90KSwgXFxpZF97Ml5rfSkiLDJdXQ==
        \[\begin{tikzcd}[ampersand replacement=\&]
                {X_D} \&\& {Y_D} \\
                \\
                {(X_D^\Sigma \times 2^k)} \&\& {(Y_D^\Sigma \times 2^k)}
                \arrow["{\phi_D}", from=1-1, to=1-3]
                \arrow["{c_{X_D}}"', from=1-1, to=3-1]
                \arrow["{c_{Y_D}}", from=1-3, to=3-3]
                \arrow["{ (\phi_D \circ (\cdot), \id_{2^k})}"', from=3-1, to=3-3]
            \end{tikzcd}\]
    \end{minipage}



    \begin{align*}
        \text{LHS Commutes} & \iff (\phi_I[\cdot], \phi_D[\cdot]) \circ c_{X_I} = c_{Y_I} \circ \phi_I                               \\
                            & \iff \forall_{x_I}^{X_I} (\phi_I[\cdot], \phi_D[\cdot]) \circ c_{X_I}(x_I) = c_{Y_I} \circ \phi_I(x_I) \\
        % & \iff \forall_{x_I}^{X_I}
    \end{align*}


    \begin{align*}
        \text{RHS commutes} & \iff\Big(\phi_D \circ (\cdot), \id_{2^k}\Big)\circ c_{X_D} = c_{Y_D} \circ \phi_D                                \\
                            & \iff \forall_{x_D}^{X_D} \Big(\phi_D \circ (\cdot), \id_{2^k}\Big)\circ c_{X_D}(x_D) = c_{Y_D} \circ \phi_D(x_D)
    \end{align*}


    So to summarize, the initial part looks like (for all $\sigma : \Sigma$):
    \begin{figure}[H]
        \includesvg[inkscapelatex=false, scale=0.4]{images/ldba-hom-xI.svg}
        \centering
        \caption{Initial part of the homomorphism}
    \end{figure}

    For the final part:
    \begin{figure}[H]
        \includesvg[inkscapelatex=false, scale=0.7]{images/ldba-hom-xD.svg}
        \centering
        \caption{Deterministic part of the homomorphism}
    \end{figure}
\end{resB}





\newpage
\section{Coalgebraic Behavioral Metrics}
This section summarizes

\todo[inline]{Reference coalgebraic behavioral metrics paper}

The goal of this section is to define a quantitative notion of bisimulation, namely a bisimulation (pseudo) metric for coalgebras.

The main idea: If you have some pseudometric $d^X$ on $X$, and also a way to lift $d^X$ to a pseudometric $d^{FX}$ on $FX$, then

Essentially the goal is to lift a metric on $X^2$ to a metric on $F(X^2)$.


First some review some definitions on pseudometrics, and some notation/assumptions.

Then we cover some category theory definitions that we will use.

\subsection{Metric review and notation}
\begin{defB}{Pseudometric Space}
    A pseudometic space is a pair $(X, d^X)$, for a set set $X$ and function $d^X : X \times X \to [0,T]_\R$ for some $T\in \R$, subject to: ($\forall_{x,y,z}^X$)
    \begin{itemize}
        \item (Symmetry) $d^X(x,y) = d^X(y,x)$
        \item (Triangle Inequality) $d^X(x,z) \le d^X(x,y) + d^X(y,z)$
        \item (Reflexivity) $d^X(x,x) = 0$
    \end{itemize}

    \textbf{(Note we are assuming a pseudometric is bounded by some fixed $T \in \R$)}
\end{defB}

\begin{defB}{Metric}
    A pseudometric space $(X,d^X)$ is a metric space iff it satisfies Definiteness, i.e.
    \begin{equation}
        \forall_{x,y}^X d^X(x,y) \neq 0 \implies d(x,y) > 0
    \end{equation}
    i.e. equivalently (under the pseudometric axioms):
    \begin{equation}
        d^X(x,y) = 0 \iff x = y
    \end{equation}

\end{defB}

\begin{defb}{Notation $T$, $F$} For the rest of this section, let $T \in \R$ be some fixed real number, and $F : \catset \to \catset$ be some endofunctor on $\catset$.
\end{defb}
\subsection{Categorical Prerequisites}

\begin{defB}[PMet]{$\catpmet$}
    Define $\catpmet$ as the category with:
    \begin{itemize}
        \item $\Obj (\catpmet) \eqdef $ all pseudo-metric spaces $(X, d^X)$.
        \item $\Hom (\catpmet) \eqdef$ all nonexpansive functions, i.e. $f: (X, d^X) \to (Y, d^Y)$ where $d^Y \circ (f, f) \le d^X$.

              (here $\le$ between functions is the pointwise order, i.e. $f \le g$ iff $f(x) \le g(x)$ for all $x \in X$)

        \item $\circ$ is the composition of functions.
        \item $\id_X$ is the identity function on $X$.
    \end{itemize}

    Note that this is indeed a category, as nonexpansive functions are closed under composition, and the identity function is nonexpansive.
\end{defB}



\begin{defB}[lift-functor-pmet]{Lifting of endofunctor from $\catset$ to $\catpmet$}
    Let $F : \catset \to \catset$. A lifting of $F$ to $\catpmet$ is a functor $\bar F$ such that TFDC:
    % https://q.uiver.app/#q=WzAsNCxbMCwxLCJcXGNhdHNldCJdLFswLDAsIlxcY2F0cG1ldCJdLFsxLDAsIlxcY2F0cG1ldCJdLFsxLDEsIlxcY2F0c2V0Il0sWzEsMCwiXFxmb3JnZXQiLDJdLFswLDMsIkYiLDJdLFsxLDIsIlxcYmFyIEYiXSxbMiwzLCJcXGZvcmdldCJdXQ==
    \[\begin{tikzcd}[ampersand replacement=\&]
            \catpmet \& \catpmet \\
            \catset \& \catset
            \arrow["{\bar F}", from=1-1, to=1-2]
            \arrow["\forget"', from=1-1, to=2-1]
            \arrow["\forget", from=1-2, to=2-2]
            \arrow["F"', from=2-1, to=2-2]
        \end{tikzcd}\]
    (in the category of small categories)


    Given a pseudometric space $(X,d)$, define the metric $d^F$ as the output $(FX, d^F) = \bar F(X,d)$.
\end{defB}


The general idea is this:
\begin{expb}{General idea for Coalgebraic Behavioral Metrics}

    Given an $F$-coalgebra $c : X \to FX$, we want to turn this into a
    $\bar F$-coalgebra $((X, d_C),  c)$ such that $c : (X, d_C) \to (FX, d_C^F)$ is nonexpansive ($c$ is a $\catpmet$ morphism).
    Note that $d \mapsto d^F \circ (c,c)$ is monotone. By Knaster Tarski, this has a fixed point. Let $d_C$ be this fixed point:
    \begin{align}
        d_C \eqdef \inf \{ d : \texttt{pseudometric}( : X^2 \to [0,T]_\R) | d_F \circ (c,c) \le d \}
    \end{align}
    (here inf is taken over the poset of pseudometrics)


    Since $d_C = d_C^F \circ (c,c)$, we have that $c$ is an isometry
\end{expb}

An important result is that $\bar F$ on objects is monotone:
\begin{resB}{$\bar F$ is monotone}
    \begin{enumerate}
        \item
              Let $F: \catset \to \catset$.
        \item
              Let $\bar F : \catpmet \to \catpmet$ be as in \refdef{lift-functor-pmet}.
        \item
              Let $(X,d_0), (X,d_1)$ be pseudometrics.

    \end{enumerate}
    Then:
    \begin{align*}
        d_0^F \le d_1^F
    \end{align*}

    \dline

    \begin{proof}
        Recall first that from \refdef{PMet}, morphisms only consist of non expansive functions.

        Thus $d_0 \le d_1 \iff $ there is some $f : \Hom^{\catpmet}((X,d_0), (X,d_1))$ such that $U f = \id_X$.


        Now assume $d_0 \le d_1$. Then take $f$ like above.

        Then $\bar F f : (FX, d_0^F) \to (FX, d_1^F)$ must exist by functorality.

        Furthermore, because $\bar F$ is a lifting, $U \bar F f = \id_X$. Thus $f_0^F \le f_1^F$.
    \end{proof}
\end{resB}

\subsection{Evaluations}
Now consider coalgebras on $\catset$.


\newcommand{\eval}[0]{\texttt{eval}}


We can think of $\catset / [0,T]_\R$ $\left(\text{recall }\refdef{slice-cat}\right)$ as having objects that have a sort of value function $V_X : X \to [0,T]_\R$, which assings a value to each state.

Note that also any pseudometric space $(X,d^X)$ can be thought of as an object in $\catset / [0,T]_\R$, namely $d^X : X\times X \to [0,T]_\R$.

\begin{defB}{Evaluation Function, Evaluation Functor}
    Given a fixed $\catset$ endofunctor $F$. We want to find a way to lift $F$ to $\catset / [0,T]_\R$:
    \begin{defb}{Evaluation Function $\eval$}
        An evaluation function is any $F$ algebra with domain $[0,T]_\R$, i.e. any map:
        \begin{equation}
            \eval : F([0,T]_\R) \to [0,T]_\R
        \end{equation}
    \end{defb}

    \begin{defb}{Evaluation Functor $\tilde F$}
        Given a fixed $\eval$, $\tilde F$ is an endofunctor on $\catset / [0,T]_\R$ that takes:
        \begin{itemize}
            \item On objects:  $V_X \mapsto \eval \circ FV_X$
            \item On morphisms: $f \mapsto F f$
        \end{itemize}

        This means you can lift any diagram in the slice category. Pictorially, this means that you can apply $\tilde F$ everywhere to a diagram:
        % https://q.uiver.app/#q=WzAsOCxbMCwwLCJYIl0sWzEsMywiWzAsVF1fXFxSIl0sWzUsMCwiRlgiXSxbNiwxLCJGWzAsVF1fXFxSIl0sWzYsMywiWzAsVF1fXFxSIl0sWzcsMCwiRlkiXSxbMiwwLCJZIl0sWzMsMSwiIFxcaHNwYWNlezJlbX1cXG1hcHN0byJdLFswLDEsIlZfWCIsMix7ImN1cnZlIjoyfV0sWzIsMywiVl9YIiwyLHsic3R5bGUiOnsiYm9keSI6eyJuYW1lIjoiZGFzaGVkIn19fV0sWzIsNCwiXFxlcWRlZiBcXHRpbGRlIEYgVl9YIiwyLHsiY3VydmUiOjJ9XSxbMyw0LCJcXGV2YWwiLDEseyJzdHlsZSI6eyJib2R5Ijp7Im5hbWUiOiJkYXNoZWQifX19XSxbMiw1LCJcXHRpbGRlIEZmXFxlcWRlZiBGZiJdLFs1LDMsIlZfWSIsMCx7InN0eWxlIjp7ImJvZHkiOnsibmFtZSI6ImRhc2hlZCJ9fX1dLFs1LDQsIlxcZXFkZWYgXFx0aWxkZSBGIFZfWSIsMCx7ImN1cnZlIjotMn1dLFswLDYsImYiXSxbNiwxLCJWX1kiLDAseyJjdXJ2ZSI6LTJ9XV0=
        \[\begin{tikzcd}[ampersand replacement=\&]
                X \&\& Y \&\&\& FX \&\& FY \\
                \&\&\& { \hspace{2em}\mapsto} \&\&\& {F[0,T]_\R} \\
                \\
                \& {[0,T]_\R} \&\&\&\&\& {[0,T]_\R}
                \arrow["f", from=1-1, to=1-3]
                \arrow["{V_X}"', curve={height=12pt}, from=1-1, to=4-2]
                \arrow["{V_Y}", curve={height=-12pt}, from=1-3, to=4-2]
                \arrow["{\tilde Ff\eqdef Ff}", from=1-6, to=1-8]
                \arrow["{V_X}"', dashed, from=1-6, to=2-7]
                \arrow["{\eqdef \tilde F V_X}"', curve={height=12pt}, from=1-6, to=4-7]
                \arrow["{V_Y}", dashed, from=1-8, to=2-7]
                \arrow["{\eqdef \tilde F V_Y}", curve={height=-12pt}, from=1-8, to=4-7]
                \arrow["\eval"{description}, dashed, from=2-7, to=4-7]
            \end{tikzcd}\]

    \end{defb}

\end{defB}

\subsection{Kantorovic Metric}

\begin{defB}{Kantorovic Metric}
    Given $F$, $\eval$, then for any pseudometric $(X,d)$ define $d_K$:
    \begin{align}
        d_{K(X,d^X)}(t_1, t_2) & \eqdef \sup \bigg \{ d^{[0,T]_\R} ( \tilde F f(t_1), \tilde f(t_2)) \bigg | f  : \Hom^{\catpmet}((X,d), ([0,T]_\R, d^{[0,T]_\R})) \bigg\} \\
                               & = \sup \Bigg \{ \Big | \tilde F V_X(t_1) - \tilde F V_X(t_2) \Big | \bigg | V_X:  X \to [0,T]_\R, \text{ $V_X$ non expansive}\bigg \}
    \end{align}
    Pictorially, it is $\sup$ over all non expansive functions $V_X : (X, d^X) \to ([0,T]_\R, d^{[0,T]_\R})$ such that TFDC:
    % https://q.uiver.app/#q=WzAsNCxbMCwwLCJGWCBcXHRpbWVzIEZYIl0sWzAsMiwiWzAsVF1fXFxSIl0sWzMsMCwiRlswLFRdX1xcUiBcXHRpbWVzIEZbMCxUXV9cXFIiXSxbMywyLCJbMCxUXV9cXFIgXFx0aW1lcyBbMCxUXV9cXFIiXSxbMCwxLCJkX0siLDIseyJzdHlsZSI6eyJib2R5Ijp7Im5hbWUiOiJkYXNoZWQifX19XSxbMCwyLCIgXFxzdXBfe1ZfWCB9IEZWX1hcXHRpbWVzIEZWX1giLDAseyJzdHlsZSI6eyJib2R5Ijp7Im5hbWUiOiJzcXVpZ2dseSJ9fX1dLFsyLDMsIlxcZXZhbCBcXHRpbWVzIFxcZXZhbCJdLFszLDEsImRee1swLFRdX1xcUn0iXV0=
    \[\begin{tikzcd}[ampersand replacement=\&]
            {FX \times FX} \&\&\& {F[0,T]_\R \times F[0,T]_\R} \\
            \\
            {[0,T]_\R} \&\&\& {[0,T]_\R \times [0,T]_\R}
            \arrow["{ \sup_{V_X } FV_X\times FV_X}", squiggly, from=1-1, to=1-4]
            \arrow["{d_K}"', dashed, from=1-1, to=3-1]
            \arrow["{\eval \times \eval}", from=1-4, to=3-4]
            \arrow["{d^{[0,T]_\R}}", from=3-4, to=3-1]
        \end{tikzcd}\]

    where $d^{[0,T]_\R}$ is $L1$ distance on $\R$.
\end{defB}

\begin{defB}{Kantorovich Lifting/Functor}
    The Kantorovich lifting (functor) of $F: \catset \to \catset$ is $F_K : \catpmet \to \catpmet$ such that $F_K (X,d^X) = (X, d_{K(X,d^X)})$, and $F_K f = Ff$.
\end{defB}

\todo[inline]{Show more examples, in particular on $\Delta$}


\begin{resb}{Useful results}
    \begin{enumerate}
        \item If $(X,d)$ is a pseudometric, then so is $d_K$.
        \item $F_K$ for a functor $F$ preserves isometries
        \item
    \end{enumerate}
\end{resb}


\begin{resb}{Kantorovich metric for MDP (new)}
    Recall the MDP functor \refdef{mdp-functor}.

    Below represents applying $\eval \circ F V_X$ for some arbitrary $V_X : X \to \R$.

    Suppose you have an element

    $$(r_0, P_{X_0}) : FX = \R \times (\Delta(X))^A$$

    After applying $F V_X$, we get:

    $$(r_0, P^\M [x_0 \lto{a} V_X^{-1}[\{\cdot\}]]) : F[0,T]_\R$$

    Now let $\eval : F[0,T]_\R \to \R$ be the evaluation function.
    This can be arbitrary, but a natural choice is to take the current reward, and sum it with the discounted expected next value.
    To get rid of $a$ we can apply $\max$. Applying this to the above:

    $$r_0 + \gamma \max_{a : A} \sum_{t \in [0,T]_\R} t P^\M [x_0 \lto{a} V_X^{-1}[\{t\}]]$$

    In particular, if $(r_0, P_{X_0})$ is the output from state $x_0$, then we get that the above is:

    $$ = r^\M(x_0) + \gamma \max_{a : A} \ev_{x_0 \lto{a} x'} [V_X(x')] $$

    Note that this is exactly the \refdef{bellman-operator} $\T$, applied to $V_X$.

    \todo[inline]{you also need to bound reward sutiably, like instead of $\R$ do $[0, (1-\gamma) T]_\R$}

    Thus the Kantorovich metric for MDPs is:

\end{resb}



\begin{expb}{Not nice metric for functor}
    Consider the functor $F : X \to A \times X$. Suppose you wanted to find the wasserstein metric.

    Note that from \refdef{coupling}, suppose $(x_0,b_0), (x_1, b_1)$ are two elements of $FX$. Then $(x_0,b_0)$ can only be coupled iff $b_0 = b_1$.

    This is not what you would expect if $B$ had some sort of metric structure already, say $B = \R$. In fact the wasserstein metric removes all metric information (namely it assumes the discrete topology on $B$ by default).

    To fix this issue, we will introduce multifunctors, and lifting pseudometrics of mutlifunctors.

\end{expb}





\newpage
\section{More complicated results}

\subsection{Limits/Colimits, Pullbacks}

\begin{defB}[limit]{Limit}
    \todo[inline]{Fill in}

\end{defB}

\begin{defb}[weak-limit]{Weak limit}
    A weak limit is just like a limit, but the map does not need to be unique.
\end{defb}


\begin{defB}[colimit]{Colimit}
    \todo[inline]{Fill in}

\end{defB}

\begin{defB}[pullback]{Pullback}
    Let $X \lto{\phi_X} Z \lfrom{\phi_Y} Y$ be a cospan. A pullback of this cospan is the limit of the diagram.

    More concretely, it is an object $P$ with maps $p_X : P \to X$ and $p_Y : P \to Y$ such that for any other object $Q$ with maps $q_X : Q\to X$, $q_Y : Q \to Y$, there is a unique map $!:Q \to P$ TFDC:
    % https://q.uiver.app/#q=WzAsNSxbMSwxLCJQIl0sWzEsMiwiWCJdLFsyLDEsIlkiXSxbMiwyLCJaIl0sWzAsMCwiUSJdLFsxLDMsIlxccGhpX1giLDJdLFsyLDMsIlxccGhpX1oiXSxbMCwxLCJwX1giLDJdLFswLDIsInBfWSJdLFs0LDEsInFfWCIsMix7ImN1cnZlIjozfV0sWzQsMiwicV9ZIiwwLHsiY3VydmUiOi0zfV0sWzQsMCwiXFxleGlzdHMgISIsMCx7InN0eWxlIjp7ImJvZHkiOnsibmFtZSI6ImRhc2hlZCJ9fX1dXQ==
    \[\begin{tikzcd}[ampersand replacement=\&]
            Q \\
            \& P \& Y \\
            \& X \& Z
            \arrow["{\exists !}", dashed, from=1-1, to=2-2]
            \arrow["{q_Y}", curve={height=-18pt}, from=1-1, to=2-3]
            \arrow["{q_X}"', curve={height=18pt}, from=1-1, to=3-2]
            \arrow["{p_Y}", from=2-2, to=2-3]
            \arrow["{p_X}"', from=2-2, to=3-2]
            \arrow["{\phi_Z}", from=2-3, to=3-3]
            \arrow["{\phi_X}"', from=3-2, to=3-3]
        \end{tikzcd}\]
\end{defB}
\begin{defB}{Weak Pullback}
    A weak pullback is a weak limit of a cospan $X \lto{\phi_X} Z \lfrom{\phi_Y} Y$. I.e. see \refdef{pullback}, but the map  $: Q \to P$ does not have to be unique:
    % https://q.uiver.app/#q=WzAsNSxbMSwxLCJQIl0sWzEsMiwiWCJdLFsyLDEsIlkiXSxbMiwyLCJaIl0sWzAsMCwiUSJdLFsxLDMsIlxccGhpX1giLDJdLFsyLDMsIlxccGhpX1oiXSxbMCwxLCJwX1giLDJdLFswLDIsInBfWSJdLFs0LDEsInFfWCIsMix7ImN1cnZlIjozfV0sWzQsMiwicV9ZIiwwLHsiY3VydmUiOi0zfV0sWzQsMCwiXFxleGlzdHMiLDAseyJzdHlsZSI6eyJib2R5Ijp7Im5hbWUiOiJkYXNoZWQifX19XV0=
    \[\begin{tikzcd}[ampersand replacement=\&]
            Q \\
            \& P \& Y \\
            \& X \& Z
            \arrow["\exists", dashed, from=1-1, to=2-2]
            \arrow["{q_Y}", curve={height=-18pt}, from=1-1, to=2-3]
            \arrow["{q_X}"', curve={height=18pt}, from=1-1, to=3-2]
            \arrow["{p_Y}", from=2-2, to=2-3]
            \arrow["{p_X}"', from=2-2, to=3-2]
            \arrow["{\phi_Z}", from=2-3, to=3-3]
            \arrow["{\phi_X}"', from=3-2, to=3-3]
        \end{tikzcd}\]

\end{defB}

\subsection{Comparing Bisimulation}
First we will need the following definitions/results:

\begin{defb}{Weak pullback preserving functor}
    Let $F : C \to D$ be a functor.

    $F$ preserves weak pullbacks iff for any pullback
    % https://q.uiver.app/#q=WzAsNCxbMCwwLCJQIl0sWzAsMSwiWCJdLFsxLDAsIlkiXSxbMSwxLCJaIl0sWzIsMywiXFxwaGlfWSJdLFsxLDMsIlxccGhpX1giLDJdLFswLDEsInBfWCIsMl0sWzAsMiwicF9ZIl1d
    \[\begin{tikzcd}[ampersand replacement=\&]
            P \& Y \\
            X \& Z
            \arrow["{p_Y}", from=1-1, to=1-2]
            \arrow["{p_X}"', from=1-1, to=2-1]
            \arrow["{\phi_Y}", from=1-2, to=2-2]
            \arrow["{\phi_X}"', from=2-1, to=2-2]
        \end{tikzcd}\]

    we have that the following commutative diagram is also a pullback of cospan $FX \lto{F\phi_X} FZ \lfrom{F\phi_Y} FY$:
    % https://q.uiver.app/#q=WzAsNCxbMCwwLCJGUCJdLFswLDEsIkZYIl0sWzEsMCwiRlkiXSxbMSwxLCJGWiJdLFsyLDMsIkZcXHBoaV9ZIl0sWzEsMywiRlxccGhpX1giLDJdLFswLDEsIkZwX1giLDJdLFswLDIsIkZwX1kiXV0=
    \[\begin{tikzcd}[ampersand replacement=\&]
            FP \& FY \\
            FX \& FZ
            \arrow["{Fp_Y}", from=1-1, to=1-2]
            \arrow["{Fp_X}"', from=1-1, to=2-1]
            \arrow["{F\phi_Y}", from=1-2, to=2-2]
            \arrow["{F\phi_X}"', from=2-1, to=2-2]
        \end{tikzcd}\]

    (meaning if there is a $Q \in \Obj (D)$ with morphisms to $FX$ and $FY$, it must factor through $FP$)

\end{defb}


% \begin{resb}{Weak pullbacks on $F$-coalgebras}
%     Let $C$ be a category with pullbacks.

%     \todo[inline]{Maybe works with weak pullbacks too but not sure how yet}

%     Suppose endofunctor $F: C \to C$ preserves weak pullbacks.

%     Then the category of $F$ coalgebras has weak pullbacks.
%     \dline

%     \begin{proof}
%         Suppose you have a cospan $c_X \lto{\phi_X} c_Z \lfrom{\phi_Y} c_Y$, where $c_X : X \to FX$, $c_Y : Y \to FY$, $c_Z : Z \to FZ$.

%         Because $C$ has weak pullbacks, take $(P, p_X : P \to X, p_Y : P \to Y)$ to be the pullback of

%         $X \lto{\phi_X} Z \lfrom{\phi_Y} Y$. Using the definition of coalgebra homomorphisms, we get TFDC:
%         % https://q.uiver.app/#q=WzAsNyxbMCwwLCJYIl0sWzIsMiwiWiJdLFs0LDQsIlkiXSxbMCwyLCJGWCJdLFsyLDQsIkZaIl0sWzQsNiwiRlkiXSxbNSwwLCJQIl0sWzAsMywiY19YIiwyXSxbMSw0LCJjX1oiLDFdLFsyLDUsImNfWSIsMl0sWzMsNCwiRiBcXHBoaV9YIiwxXSxbMCwxLCJcXHBoaV9YIiwxXSxbMiwxLCJcXHBoaV9ZIiwxXSxbNSw0LCJGXFxwaGlfWSIsMV0sWzYsMCwicF9YIiwxXSxbNiwyXV0=
%         \[\begin{tikzcd}[ampersand replacement=\&]
%                 X \&\&\&\&\& P \\
%                 \\
%                 FX \&\& Z \\
%                 \\
%                 \&\& FZ \&\& Y \\
%                 \\
%                 \&\&\&\& FY
%                 \arrow["{c_X}"', from=1-1, to=3-1]
%                 \arrow["{\phi_X}"{description}, from=1-1, to=3-3]
%                 \arrow["{p_X}"{description}, from=1-6, to=1-1]
%                 \arrow[from=1-6, to=5-5]
%                 \arrow["{F \phi_X}"{description}, from=3-1, to=5-3]
%                 \arrow["{c_Z}"{description}, from=3-3, to=5-3]
%                 \arrow["{\phi_Y}"{description}, from=5-5, to=3-3]
%                 \arrow["{c_Y}"', from=5-5, to=7-5]
%                 \arrow["{F\phi_Y}"{description}, from=7-5, to=5-3]
%             \end{tikzcd}\]

%         Then because $F$ preserves weak pullbacks, we can add the pullback $(FP, Fp_X : FP \to FX, Fp_Y : FP \to FY)$. Then notice that $c_X \circ p_X : P \to FX$, and $c_Y \circ p_Y : P \to FY$.

%         So becasue $FP$ is a pullback, there must be a map $f : P \to FP$. As a diagram, this means TFDC:
%         % https://q.uiver.app/#q=WzAsOCxbMCwwLCJYIl0sWzIsMiwiWiJdLFs0LDQsIlkiXSxbMCwyLCJGWCJdLFsyLDQsIkZaIl0sWzQsNiwiRlkiXSxbNSwyLCJGUCJdLFs1LDAsIlAiXSxbMCwzLCJjX1giLDJdLFsxLDQsImNfWiIsMV0sWzIsNSwiY19ZIiwyXSxbMyw0LCJGIFxccGhpX1giLDFdLFswLDEsIlxccGhpX1giLDFdLFsyLDEsIlxccGhpX1kiLDFdLFs1LDQsIkZcXHBoaV9ZIiwxXSxbNiwzLCJGcF9YIiwxXSxbNiw1LCJGcF9ZIiwxXSxbNywwLCJwX1giLDFdLFs3LDJdLFs3LDYsIlxcZXhpc3RzICBjX1AiLDEseyJzdHlsZSI6eyJib2R5Ijp7Im5hbWUiOiJkYXNoZWQifX19XV0=
%         \[\begin{tikzcd}[ampersand replacement=\&]
%                 X \&\&\&\&\& P \\
%                 \\
%                 FX \&\& Z \&\&\& FP \\
%                 \\
%                 \&\& FZ \&\& Y \\
%                 \\
%                 \&\&\&\& FY
%                 \arrow["{c_X}"', from=1-1, to=3-1]
%                 \arrow["{\phi_X}"{description}, from=1-1, to=3-3]
%                 \arrow["{p_X}"{description}, from=1-6, to=1-1]
%                 \arrow["{\exists  c_P}"{description}, dashed, from=1-6, to=3-6]
%                 \arrow[from=1-6, to=5-5]
%                 \arrow["{F \phi_X}"{description}, from=3-1, to=5-3]
%                 \arrow["{c_Z}"{description}, from=3-3, to=5-3]
%                 \arrow["{Fp_X}"{description}, from=3-6, to=3-1]
%                 \arrow["{Fp_Y}"{description}, from=3-6, to=7-5]
%                 \arrow["{\phi_Y}"{description}, from=5-5, to=3-3]
%                 \arrow["{c_Y}"', from=5-5, to=7-5]
%                 \arrow["{F\phi_Y}"{description}, from=7-5, to=5-3]
%             \end{tikzcd}\]
%         Thus in the category of coalgebras:
%         % https://q.uiver.app/#q=WzAsNCxbMywwLCJjX1AiXSxbMCwwLCJjX1giXSxbMSwxLCJjX1oiXSxbMiwyLCJjX1kiXSxbMCwxLCJwX1giLDJdLFswLDMsInBfWSJdLFsxLDIsIlxccGhpX1giLDJdLFszLDIsIlxccGhpX1kiXV0=
%         \[\begin{tikzcd}[ampersand replacement=\&]
%                 {c_X} \&\&\& {c_P} \\
%                 \& {c_Z} \\
%                 \&\& {c_Y}
%                 \arrow["{\phi_X}"', from=1-1, to=2-2]
%                 \arrow["{p_X}"', from=1-4, to=1-1]
%                 \arrow["{p_Y}", from=1-4, to=3-3]
%                 \arrow["{\phi_Y}", from=3-3, to=2-2]
%             \end{tikzcd}\]
%         Now take any other coalgebra $c_Q : Q \to FQ$, with homomorphisms $q_X : c_Q \to c_X$, and $q_Y : c_Q \to c_Y$.

%         Because $P$ is a pullback, there must exist some $g : Q \to P$ that commutes with the pullback square $(P,X,Y,Z)$.

%         Likewise from functorality, $Fg: FQ \to FP$ must commute with the pullback square $(FP, FX, FY, FZ)$.



%         Thus the following diagram almost commutes:
%         % https://q.uiver.app/#q=WzAsMTAsWzAsMiwiWCJdLFsyLDQsIloiXSxbNCw2LCJZIl0sWzAsNCwiRlgiXSxbMiw2LCJGWiJdLFs0LDgsIkZZIl0sWzgsMCwiUSJdLFs4LDIsIkZRIl0sWzUsNCwiRlAiXSxbNSwyLCJQIl0sWzAsMywiY19YIiwyXSxbMSw0LCJjX1oiLDFdLFsyLDUsImNfWSIsMl0sWzMsNCwiRiBcXHBoaV9YIiwxXSxbMCwxLCJcXHBoaV9YIiwxXSxbMiwxLCJcXHBoaV9ZIiwxXSxbNSw0LCJGXFxwaGlfWSIsMV0sWzYsMCwicV9YIiwxLHsibGFiZWxfcG9zaXRpb24iOjIwLCJjdXJ2ZSI6Mn1dLFs2LDIsInFfWSIsMSx7ImxhYmVsX3Bvc2l0aW9uIjoyMCwiY3VydmUiOi00fV0sWzcsMywiRnFfWCIsMSx7ImxhYmVsX3Bvc2l0aW9uIjoyMCwiY3VydmUiOjJ9XSxbNiw3LCJjX1EiXSxbOCwzLCJGcF9YIiwxXSxbOCw1LCJGcF9ZIiwxXSxbOSwwLCJwX1giLDFdLFs5LDJdLFs5LDgsIlxcZXhpc3RzICBjX1AiLDEseyJzdHlsZSI6eyJib2R5Ijp7Im5hbWUiOiJkYXNoZWQifX19XSxbNiw5LCJcXGV4aXN0cyBnIiwxLHsic3R5bGUiOnsiYm9keSI6eyJuYW1lIjoiZGFzaGVkIn19fV0sWzcsNSwiRnFfWSIsMSx7ImN1cnZlIjotNH1dLFs3LDgsIkZnIiwxXV0=
%         \[\begin{tikzcd}[ampersand replacement=\&]
%             \&\&\&\&\&\&\&\& Q \\
%             \\
%             X \&\&\&\&\& P \&\&\& FQ \\
%             \\
%             FX \&\& Z \&\&\& FP \\
%             \\
%             \&\& FZ \&\& Y \\
%             \\
%             \&\&\&\& FY
%             \arrow["{q_X}"{description, pos=0.2}, curve={height=12pt}, from=1-9, to=3-1]
%             \arrow["{\exists g}"{description}, dashed, from=1-9, to=3-6]
%             \arrow["{c_Q}", from=1-9, to=3-9]
%             \arrow["{q_Y}"{description, pos=0.2}, curve={height=-24pt}, from=1-9, to=7-5]
%             \arrow["{c_X}"', from=3-1, to=5-1]
%             \arrow["{\phi_X}"{description}, from=3-1, to=5-3]
%             \arrow["{p_X}"{description}, from=3-6, to=3-1]
%             \arrow["{\exists  c_P}"{description}, dashed, from=3-6, to=5-6]
%             \arrow[from=3-6, to=7-5]
%             \arrow["{Fq_X}"{description, pos=0.2}, curve={height=12pt}, from=3-9, to=5-1]
%             \arrow["Fg"{description}, from=3-9, to=5-6]
%             \arrow["{Fq_Y}"{description}, curve={height=-24pt}, from=3-9, to=9-5]
%                 \arrow["{F \phi_X}"{description}, from=5-1, to=7-3]
%                 \arrow["{c_Z}"{description}, from=5-3, to=7-3]
%                 \arrow["{Fp_X}"{description}, from=5-6, to=5-1]
%                 \arrow["{Fp_Y}"{description}, from=5-6, to=9-5]
%                 \arrow["{\phi_Y}"{description}, from=7-5, to=5-3]
%                 \arrow["{c_Y}"', from=7-5, to=9-5]
%                 \arrow["{F\phi_Y}"{description}, from=9-5, to=7-3]
%             \end{tikzcd}\]

%             All that remains to show is that TFDC:
%                     % https://q.uiver.app/#q=WzAsNCxbMCwwLCJQIl0sWzAsMSwiRlAiXSxbMSwwLCJRIl0sWzEsMSwiRlEiXSxbMCwxLCJjX1AiLDJdLFsyLDMsImNfUSJdLFsyLDAsImciLDJdLFszLDEsIkZnIl1d
%                     \[\begin{tikzcd}[ampersand replacement=\&]
%                             P \& Q \\
%                             FP \& FQ
%                             \arrow["{c_P}"', from=1-1, to=2-1]
%                             \arrow["g"', from=1-2, to=1-1]
%                             \arrow["{c_Q}", from=1-2, to=2-2]
%                             \arrow["Fg", from=2-2, to=2-1]
%                         \end{tikzcd}\]
%         Note that to show this fact, because $C$ has pullbacks, note that $Q$ is a cone over $FX \lto{F\phi_X} FZ \lfrom{F\phi_Y} $



%         Thus in the category of $F$ coalgebras, we have:
%         % https://q.uiver.app/#q=WzAsNSxbMywxLCJjX1AiXSxbMCwxLCJjX1giXSxbMSwyLCJjX1oiXSxbMiwzLCJjX1kiXSxbNSwwLCJjX1EiXSxbMCwxLCJwX1giLDJdLFswLDMsInBfWSJdLFsxLDIsIlxccGhpX1giLDJdLFszLDIsIlxccGhpX1kiXSxbNCwxLCJxX1giLDIseyJjdXJ2ZSI6Mn1dLFs0LDMsInFfWSIsMCx7ImN1cnZlIjotM31dLFs0LDAsIlxcZXhpc3RzIGciLDEseyJzdHlsZSI6eyJib2R5Ijp7Im5hbWUiOiJkb3R0ZWQifX19XV0=
%         \[\begin{tikzcd}[ampersand replacement=\&]
%                 \&\&\&\&\& {c_Q} \\
%                 {c_X} \&\&\& {c_P} \\
%                 \& {c_Z} \\
%                 \&\& {c_Y}
%                 \arrow["{q_X}"', curve={height=12pt}, from=1-6, to=2-1]
%                 \arrow["{\exists g}"{description}, dotted, from=1-6, to=2-4]
%                 \arrow["{q_Y}", curve={height=-18pt}, from=1-6, to=4-3]
%                 \arrow["{\phi_X}"', from=2-1, to=3-2]
%                 \arrow["{p_X}"', from=2-4, to=2-1]
%                 \arrow["{p_Y}", from=2-4, to=4-3]
%                 \arrow["{\phi_Y}", from=4-3, to=3-2]
%             \end{tikzcd}\]

%         which shows that this category has weak pullbacks.
%     \end{proof}
% \end{resb}
\begin{resB}[AM-vs-beh-equiv]{AM Bisimulation is weaker than Behavioral Equivalence}
    Suppose $F : C \to C$ preserves weak pullbacks, and $C$ has weak pullbacks. Then for any AM bisimulation $c_X \lfrom{\phi_X} c_R \lto{\phi_Y} c_Y$ there is a behavioral equivalence.


    \dline

    % \begin{proof} \hfill

    We will use the following:

    % https://q.uiver.app/#q=WzAsNCxbMCwxLCJjX1giXSxbMSwwLCJjX1kiXSxbMSwxLCIqIl0sWzAsMCwiXFxleGlzdHMgY19SIl0sWzAsMl0sWzEsMl0sWzMsMF0sWzMsMV1d
    \[\begin{tikzcd}[ampersand replacement=\&]
            {\exists c_R} \& {c_Y} \\
            {c_X} \& {*}
            \arrow[from=1-1, to=1-2]
            \arrow[from=1-1, to=2-1]
            \arrow[from=1-2, to=2-2]
            \arrow[from=2-1, to=2-2]
        \end{tikzcd}\]


    So basically we want to show taht given a

    % \end{proof}

\end{resB}


\subsection{Existance of final coalgebra}
(A slightly more involved result shows how to build a final coalgebra)


\begin{defB}{(Ordinal Indexed) Chain, Cochain (new?)}
    Given an ordinal $\alpha$, an $\alpha$-chain is a diagram of the form:
    % https://q.uiver.app/#q=WzAsOCxbMCwwLCJBXzAiXSxbMSwwLCJBXzEiXSxbMiwwLCJBXzIiXSxbMywwLCJcXGNkb3RzIl0sWzQsMCwiQV9cXG9tZWdhIFxcY29uZyBcXG1hdGhybXtjb2xpbX0oQV8wXFx0b1xcY2RvdHMpIl0sWzUsMCwiQV97XFxvbWVnYSArIDF9Il0sWzYsMCwiXFxjZG90cyJdLFs3LDAsIkFfXFxhbHBoYSJdLFswLDFdLFsxLDJdLFsyLDNdLFszLDRdLFs0LDVdLFs1LDZdLFs2LDddXQ==
    \[\begin{tikzcd}[ampersand replacement=\&]
            {A_0} \& {A_1} \& {A_2} \& \cdots \& {A_\omega \cong \mathrm{colim}(A_0\to\cdots)} \& {A_{\omega + 1}} \& \cdots \& {A_\alpha}
            \arrow[from=1-1, to=1-2]
            \arrow[from=1-2, to=1-3]
            \arrow[from=1-3, to=1-4]
            \arrow[from=1-4, to=1-5]
            \arrow[from=1-5, to=1-6]
            \arrow[from=1-6, to=1-7]
            \arrow[from=1-7, to=1-8]
        \end{tikzcd}\]

    Likewise a $\alpha$-cochain is the dual, i.e. a diagram of the form:
    % https://q.uiver.app/#q=WzAsOCxbMCwwLCJBXzAiXSxbMSwwLCJBXzEiXSxbMiwwLCJBXzIiXSxbMywwLCJcXGNkb3RzIl0sWzQsMCwiQV9cXG9tZWdhIFxcY29uZyBcXG1hdGhybXtsaW19KEFfMFxcbGVmdGFycm93XFxjZG90cykiXSxbNSwwLCJBX3tcXG9tZWdhICsgMX0iXSxbNiwwLCJcXGNkb3RzIl0sWzcsMCwiQV9cXGFscGhhIl0sWzEsMF0sWzIsMV0sWzMsMl0sWzQsM10sWzUsNF0sWzYsNV0sWzcsNl1d
    \[\begin{tikzcd}[ampersand replacement=\&]
            {A_0} \& {A_1} \& {A_2} \& \cdots \& {A_\omega \cong \mathrm{lim}(A_0\leftarrow\cdots)} \& {A_{\omega + 1}} \& \cdots \& {A_\alpha}
            \arrow[from=1-2, to=1-1]
            \arrow[from=1-3, to=1-2]
            \arrow[from=1-4, to=1-3]
            \arrow[from=1-5, to=1-4]
            \arrow[from=1-6, to=1-5]
            \arrow[from=1-7, to=1-6]
            \arrow[from=1-8, to=1-7]
        \end{tikzcd}\]
\end{defB}

%(https://drops.dagstuhl.de/storage/00lipics/lipics-vol029-fsttcs2014/LIPIcs.FSTTCS.2014.403/LIPIcs.FSTTCS.2014.403.pdf)
\begin{resB}{Final coalgebra exists (Pohlova '73, Adamek '74)}
    Let $C$ be a cateogory with terminal object $1$.

    Consider the following $\omega$-cochain:

    % https://q.uiver.app/#q=WzAsNSxbMCwyLCIxIl0sWzEsMiwiRigxKSJdLFsyLDIsIkZeMigxKSJdLFsyLDAsIkxfRj0gXFxsaW0oMSBcXGxlZnRhcnJvdyBGKDEpIFxcbGVmdGFycm93IFxcY2RvdHMpIl0sWzMsMiwiXFxjZG90cyJdLFsxLDAsIlxcZXhpc3RzICFmIiwwLHsic3R5bGUiOnsiYm9keSI6eyJuYW1lIjoiZGFzaGVkIn19fV0sWzIsMSwiRihmKSJdLFszLDAsIlxcZXhpc3RzICFsXzAiLDEseyJzdHlsZSI6eyJib2R5Ijp7Im5hbWUiOiJkYXNoZWQifX19XSxbMywxLCJsXzEiLDFdLFszLDIsImxfMiIsMV0sWzQsMl0sWzMsNF1d
    \[\begin{tikzcd}[ampersand replacement=\&]
            \&\& {L_F= \lim(1 \leftarrow F(1) \leftarrow \cdots)} \\
            \\
            1 \& {F(1)} \& {F^2(1)} \& \cdots
            \arrow["{\exists !l_0}"{description}, dashed, from=1-3, to=3-1]
            \arrow["{l_1}"{description}, from=1-3, to=3-2]
            \arrow["{l_2}"{description}, from=1-3, to=3-3]
            \arrow[from=1-3, to=3-4]
            \arrow["{\exists !f}", dashed, from=3-2, to=3-1]
            \arrow["{F(f)}", from=3-3, to=3-2]
            \arrow[from=3-4, to=3-3]
        \end{tikzcd}\]
    Suppose that $L_F$ exists, and that $F$ preserves this limit. Then:
    \begin{itemize}
        \item There exists a coalgebra $: L_F \to F(L_F)$
        \item This coalgbra is final
    \end{itemize}

    \dline
    \begin{proof}

        Now consider shifted chain:
        % https://q.uiver.app/#q=WzAsNCxbMCwyLCJGKDEpIl0sWzEsMiwiRl4yKDEpIl0sWzEsMCwiTF9GXntcXG1hdGhybXtTSH19PSBcXGxpbSggRigxKSBcXGxlZnRhcnJvdyBcXGNkb3RzKSJdLFsyLDIsIlxcY2RvdHMiXSxbMSwwLCJGKGYpIl0sWzIsMCwibF8xIiwxXSxbMiwxLCJsXzIiLDFdLFszLDFdLFsyLDNdXQ==
        \[\begin{tikzcd}[ampersand replacement=\&]
                \& {L_F^{\mathrm{SH}}= \lim( F(1) \leftarrow \cdots)} \\
                \\
                {F(1)} \& {F^2(1)} \& \cdots
                \arrow["{l_1}"{description}, from=1-2, to=3-1]
                \arrow["{l_2}"{description}, from=1-2, to=3-2]
                \arrow[from=1-2, to=3-3]
                \arrow["{F(f)}", from=3-2, to=3-1]
                \arrow[from=3-3, to=3-2]
            \end{tikzcd}\]
        Since $F$ preserves this limit, we have that $F(L_F)$ is a limit of the shifted chain and $l_1, \ldots$. Thus we have two limts so they must be isomorphic: $\mathrm{pres}_F : F(L_F) \cong L_F^{\mathrm{SH}}$.

        Note $L_F^{\mathrm{SH}}$ is also a cone over $1 \lfrom{f} F(1) \leftarrow \cdots$ (take $!: L_F^{\mathrm{SH}} \to 1$)

        ($L_F$ is a limit) $\implies  \exists ! : L_F^{\mathrm{SH}} \to L_F$.

        Note $L_F$ is a cone over $F(1) \lfrom{F(f)} F^2(1) \leftarrow \cdots$

        ($L_F^{\mathrm{SH}}$ is a limit ) $\implies \exists ! : L_F \to L_F^{\mathrm{SH}}$.

        Thus $L_F \cong L_F^{\mathrm{SH}}$.

        Now take the coalgebra $L_F \to L_F^{\mathrm{SH}} \lto{\mathrm{pres}_F} F(L_F)$.

    \end{proof}


\end{resB}


\end{document}
